\documentclass[UTF8,a4paper,11pt,openany,fontset=fandol]{ctexbook}
\usepackage[a4paper,margin=24mm,headheight=15pt]{geometry}
\usepackage{amsmath,amssymb,bm,mathtools,booktabs,longtable,tabularx,array,graphicx}
\usepackage{enumitem,xcolor,fancyhdr,needspace}
\usepackage[numbers,sort&compress]{natbib}
\usepackage{hyperref,bookmark}
\usepackage{tikz}
\usetikzlibrary{arrows.meta,positioning}
\setCJKmainfont{FandolSong-Regular.otf}[BoldFont=FandolSong-Bold.otf,ItalicFont=FandolKai-Regular.otf]
\setCJKsansfont{FandolHei-Regular.otf}[BoldFont=FandolHei-Bold.otf]
\setCJKmonofont{FandolFang-Regular.otf}
\setmainfont{texgyretermes-regular.otf}[BoldFont=texgyretermes-bold.otf,ItalicFont=texgyretermes-italic.otf,BoldItalicFont=texgyretermes-bolditalic.otf]\setsansfont{texgyreheros-regular.otf}[BoldFont=texgyreheros-bold.otf]
\definecolor{ink}{HTML}{193D55}
\ctexset{chapter={format=\sffamily\huge\bfseries\color{ink},beforeskip=0pt,afterskip=20pt},section={format=\sffamily\Large\bfseries\color{ink}},subsection={format=\sffamily\large\bfseries},subsubsection={format=\sffamily\normalsize\bfseries}}
\hypersetup{colorlinks=true,linkcolor=ink,citecolor=ink,urlcolor=ink,bookmarksnumbered=true,bookmarksdepth=2,pdftitle={Hall MHD与双流体磁重联：模型推导、数值方法与文献调研}}
\setlength{\emergencystretch}{4em}\setlength{\parskip}{3pt}\linespread{1.12}\raggedbottom
\setlist{nosep,leftmargin=2em}\setcounter{tocdepth}{1}\setcounter{secnumdepth}{2}
\pagestyle{fancy}\fancyhf{}\fancyhead[L]{\small\nouppercase{\leftmark}}\fancyfoot[C]{\thepage}
\newcommand{\vct}[1]{\bm{#1}}\newcommand{\dd}{\mathrm d}\newcommand{\VA}{V_{\!A}}\newcommand{\J}{\vct J}\newcommand{\B}{\vct B}\newcommand{\E}{\vct E}\newcommand{\vel}{\vct u}\newcommand{\grad}{\nabla}\newcommand{\R}{\mathcal R}
\newcommand{\note}[1]{\par\smallskip\noindent\textsf{\textbf{学习提示：}}#1\par\smallskip}
\newcommand{\assume}[1]{\par\smallskip\noindent\textsf{\textbf{假设与边界：}}#1\par\smallskip}
\newcolumntype{P}[1]{>{\raggedright\arraybackslash}p{#1}}
\numberwithin{equation}{chapter}

\let\cleardoublepage\clearpage

\tracinglostchars=3

\hypersetup{pdftitle={双流体与流体闭合怎样解释磁重联中的离子电子加热}}
\newcommand{\loc}[2]{\par\noindent\textsf{原文定位：}\texttt{#1}，#2。\par}
\newcommand{\eV}{\,\mathrm{eV}}
\begin{document}
\frontmatter
\begin{titlepage}
\vspace*{16mm}
{\sffamily\Large Hall MHD 与双流体磁重联调研 · 专题二\par}
\vspace{12mm}
{\sffamily\Huge\bfseries\color{ink}双流体与流体闭合\\[8pt]怎样解释离子、电子加热\par}
\vspace{12mm}
{\Large 从独立物种能量到压力张量、热流和实验检验\par}
\vfill
调研日期：2026年11月14日。\\[7pt]
三个比较问题：增加哪些方程才能得到新的加热结论；与实验对到哪一步；能独立支持、修正或限制哪一种 Hall MHD。\\[7pt]
主线为流体模型；动理学只用于辨认闭合边界，不作为双流体成果计数。\\[7pt]
本报告复用已有全文证据，并新增核读公开原文。逐篇版本、图式、访问缺口和计数见附录与本目录核读清单。没有运行文献中的大规模等离子体模拟。
\end{titlepage}
\chapter*{三个问题的回答}
\textbf{有超出特定 Hall 基线的结论，但增加自由度比模型名称更重要。}单总压或固定温度比模型不能独立预测两物种温升；独立能量方程可以给出碰撞交换与升温滞后。标量双温 Hall 仍不能自洽演化完整压力各向异性；十矩增加压力张量，使压力支撑、应变功和部分温度各向异性成为输出。热流闭合则决定这些压力和热点能否保持。已有研究显示，改变闭合能显著改变岛合并、层宽及三维不稳定性，即使重联率或全域热能看起来相近。\citep{L164,Wang2015,Ng2017,Ng2020,TenBarge2019MMS,Closure2026}

\textbf{实验支持具有分项边界。}MAST、HIT-SI提供温度结构和趋势约束，尚不足以唯一确认黏性、电阻或超电阻热源。MRX能量预算约束电功、储能及输运，不能把常见的离子:电子约$2:1$电场沉积改称不可逆加热比。MMS可直接检验各向异性电子状态方程；新增的2019与2021年证据尤其显示，近X线的温度转折、远下游的上游来源选择和垂直增温偏差必须分别讨论。它们是电子闭合的检验，数值部分含混合模型，不能冒充完整双流体分配验证。\citep{Tanabe2015MAST,Kaptanoglu2020,F674,Wetherton2019,Wetherton2021}

\textbf{现有证据主要提供有条件支持和明确限制，未建立普适加热比例。}双温慢激波支持“离子流能先转内能，电子随后通过交换和传导响应”这一可能路径，却不是无碰撞Hall扩散区的独立复现。十矩与Hall若共用输运参数、边界或校准数据，其相符不构成独立物理验证。最有价值的下一步是同初边值、同控制体的模型消融，再通过合成诊断比较温度、各向异性、热流和能量残差，而不是只比较峰值。

\textbf{阅读路线：}第1章规定证据与Hall基线；第2章推导独立物种方程和加热账本；第3章按问题串联流体研究；第4章单独审视实验；第5章给出比较判决和可否证研究；附录列逐篇证据、数值细节及待办。与旧专题的互补关系见\textsf{../离子加热问题/README.md}。
\begingroup\small\setcounter{tocdepth}{2}\tableofcontents\endgroup
\mainmatter
\chapter{先指定问题、模型和证据边界}
\section{这次扩大到哪里}
研究对象是非相对论磁重联中离子与电子的能量转化、温升和热输运。主要纳入有独立流体演化或可明确应用到流体方程的闭合：共速双温碰撞流体、双温Hall及惯性扩展、五矩、十矩、压力各向异性状态方程、Landau/陀螺流体，以及按来源分群的电子多流体。两种矩阶数与两种电磁约化可以交叉组合；“十矩”不自动意味着全Maxwell，“全Maxwell”也不自动意味着有正确热流。

离子--中性“两流体”、相对论电子--正电子两流体不作为本轮的离子--电子主线；并非它们不重要，而是物种身份与目标能量账本不同。PIC、Vlasov和混合模型只承担必要的参照或闭合观测检验，不扩展为另一份动理学综述。论文用双流体方程诊断PIC，也不等于独立推进了双流体模型。

检索使用网页检索、出版机构、arXiv、MIT与OSTI等公开来源，结合本库原文和引文追踪。关键词覆盖 two-fluid/ten-moment/twenty-moment/anisotropic pressure/heat flux/ion electron heating/reconnection，以及MAST、HIT-SI、MMS。逐次查询与访问状态保存在\path{search/}。搜索命中和题名筛查不是全文核读；本轮没有穷尽数据库、递归穷尽引文，也没有把未取得的正文计为新读。

旧库的全文记录按稳定键复用，定向复查不增加计数。新增读到的接收稿保留原文件、哈希及物理页码；正文、方法、讨论、图表及文内适用附录均读完，并直接查看关键PDF页后，才进入新核读记录。前序方法或独立补充尚未取得时明确保留缺口。报告中的“未报告”指本次所读版本没有提供足以支持该判断的信息，不意味着作者从未做过相应工作。

\section{Hall基线不能只有一个名字}
\begin{longtable}{P{16mm}P{54mm}P{78mm}}
\toprule 基线 & 必须写明的方程 & 哪类比较才有意义\\\midrule\endfirsthead
\toprule 基线 & 必须写明的方程 & 哪类比较才有意义\\\midrule\endhead
H0 & 等温或规定$p_e,p_i$的Hall & 温度不独立演化；任何温升预测先需要能量方程。\\
H1 & Hall欧姆律、总压或总能，固定物种分配 & 能比较总内能，但不能把指定的$T_i/T_e$当输出。\\
H2 & 独立标量$p_i,p_e$；各自对流；物理热源与输运 & 可比较温度、交换和损失；仍不能输出完整压力张量。Browning与HIT-SI是具体实现，不是统一标准。\\
H3 & H2加电子惯性和明确能量一致性 & 可比较电子动能储存与惯性破冻结；仍无自动产生的不可逆电子热源。增强电子质量也是参数。\\
H4 & Hall/惯性约化加各向异性或完整压力演化 & 与十矩的差别须再检查密度约束、位移电流、热流及碰撞，不能把压力自由度的成功归为“战胜所有Hall”。\\
\bottomrule\end{longtable}

这里H0--H4是本报告的比较标签，并非文献正式命名。H2和H3并非严格嵌套的现成代码：热源、壁面和电子惯性的实现可能同时变化。将两份不同几何的计算排在一张图上，只能构成跨条件启发，不能当成控制变量实验。

\section{成功证据分成四个层级}
第一层是模型能力：方程包含所需自由度，例如十矩能表示$P_{xy}$。第二层是数值可信度：给定闭合下，目标观测量随网格、步长和数值限制器收敛。第三层是模型比较：在匹配参数下，一个闭合比另一个更接近参照。第四层才是观测验证：经过测量响应、采样和不确定度处理后，与实验一致。层级之间没有自动推论。

对加热还须多一道判别：测量的温度峰、积分内能、焓流、电场沉积、模型熵产或分布函数尾分别是什么。相近重联率只说明磁通转移的一项指标相近；温度峰相近可能由不同热源和不同损失相抵得到。两者都不足以独自证明机制。

本报告将已解决的表述限定为“某一条件下某一可观测量已得到有证据的解释”。“模型表示不了”“闭合不可靠”“没有做同参数对照”“原文未取得”分别是不同缺口。读者应据此判断下一步是增添方程、改善闭合、设计对照，还是继续取得证据。

\chapter{从物种方程到可以检验的加热账本}
\section{未知量与矩的定义}
取单电荷离子$q_i=+e$、电子$q_e=-e$，$e>0$。先不施加准中性。密度$n_i,n_e$、平均速度$\vct u_i,\vct u_e$和中心压力张量$\vct P_i,\vct P_e$分别定义为
\begin{align}
n_i&=\int f_i\dd^3v,&n_e&=\int f_e\dd^3v,\\
\vct u_i&=n_i^{-1}\int\vct v f_i\dd^3v,&\vct u_e&=n_e^{-1}\int\vct v f_e\dd^3v,\\
P_{i,ab}&=m_i\int c_{i,a}c_{i,b}f_i\dd^3v,&P_{e,ab}&=m_e\int c_{e,a}c_{e,b}f_e\dd^3v,
\end{align}
其中$\vct c_i=\vct v-\vct u_i$、$\vct c_e=\vct v-\vct u_e$。下标$a,b,c$表示笛卡尔分量，重复指标求和。三阶中心矩和热流的关系为
\begin{equation}
Q_{s,abc}=m_s\int c_{s,a}c_{s,b}c_{s,c}f_s\dd^3v,
\qquad q_{s,a}=\frac12 Q_{s,abb},\quad s=i,e.
\end{equation}
这里$s$只在已分别定义两物种之后缩写相同结构。压力单位Pa，$Q$与热流均为$\mathrm{W\,m^{-2}}$，但一个是三阶张量，一个是向量，不能互换。文献对平行三阶矩与平行能流可能使用不同二分之一约定，本报告以本式为准。

定义$p_s=\operatorname{tr}\vct P_s/3$、$T_s=p_s/n_s$，温度以能量单位表示；$T_s/k_B$才是K。内能$U_s=\operatorname{tr}\vct P_s/2$，流能$K_s=m_sn_su_s^2/2$。这些定义不要求分布为Maxwell型，因此非热多束也能贡献二阶矩。\textbf{有限矩模型并不唯一确定热核心、尾部及精细分布熵。}

\section{分别列出的完整非相对论两物种方程}
无反应、无辐射及无外源时，连续性为
\begin{align}
\partial_t n_i+\nabla\cdot(n_i\vct u_i)&=0,\\
\partial_t n_e+\nabla\cdot(n_e\vct u_e)&=0.
\end{align}
动量方程为
\begin{align}
m_in_iD_i\vct u_i&=en_i(\E+\vct u_i\times\B)-\nabla\cdot\vct P_i+\vct R_i,\\
m_en_eD_e\vct u_e&=-en_e(\E+\vct u_e\times\B)-\nabla\cdot\vct P_e+\vct R_e,
\end{align}
其中$D_i=\partial_t+\vct u_i\cdot\nabla$，$D_e=\partial_t+\vct u_e\cdot\nabla$；$\vct R_s$是碰撞动量交换密度。两套内能方程是
\begin{align}
\partial_tU_i+\nabla\cdot(U_i\vct u_i+\vct q_i)&=-\vct P_i:\nabla\vct u_i+H_i,\label{eq:Ui}\\
\partial_tU_e+\nabla\cdot(U_e\vct u_e+\vct q_e)&=-\vct P_e:\nabla\vct u_e+H_e.\label{eq:Ue}
\end{align}
此处$H_s$是碰撞或明确给定的内能源密度，单位$\mathrm{W\,m^{-3}}$。两个方程中的速度不能擅自统一；电子压缩、对流和焓流通常沿电子流。

全电磁方程及约束为
\begin{align}
\partial_t\B&=-\nabla\times\E,&\epsilon_0\partial_t\E&=\mu_0^{-1}\nabla\times\B-\J,\\
\nabla\cdot\B&=0,&\epsilon_0\nabla\cdot\E&=e(n_i-n_e),\\
\J&=e(n_i\vct u_i-n_e\vct u_e).
\end{align}
五矩取每物种$n_s$、三个动量和一个总粒子能量$K_s+U_s$，假定$\vct P_s=p_s\vct I$并指定热流、碰撞；两物种共有十个流体变量，另有六个电磁演化分量及散度约束。十矩以六个对称压力分量替代标量能量，每物种共十个矩；内能是压力迹，不能再额外独立推进一个不协调的标量能量方程。\citep{Wang2015,TenBarge2019MMS}

\section{Hall约化删去什么、保留什么}
在$\lambda_D/L\ll1$、低频弱电荷分离条件下令$n_i\simeq n_e=n$，在电磁传播时间远短于演化时间的排序下忽略位移电流。于是$\J=\nabla\times\B/\mu_0$且
\begin{equation}\vct u_e=\vct u_i-\frac{\J}{en}.\end{equation}
把电子动量方程除以$-en$再代入，得
\begin{equation}
\E+\vct u_i\times\B=
\frac{\J\times\B}{en}-\frac{\nabla\cdot\vct P_e}{en}
-\frac{m_e}{e}D_e\vct u_e+\frac{\vct R_e}{en}.
\label{eq:ohm}
\end{equation}
令$m_eD_e\vct u_e$在目标尺度上可忽略，就得到无电子惯性的Hall形式；若再取标量压力，则失去非旋转对称压力支撑。电子惯性的删除必须以频率和层厚排序说明，不能仅凭总体质量比小就在电子层删除。准中性也不等于$\vct u_i=\vct u_e$。

质心速度$\vct u=(m_i\vct u_i+m_e\vct u_e)/(m_i+m_e)$下，两物种总流能严格分为
\begin{equation}
K_i+K_e=\frac{\rho u^2}{2}+\frac{m_im_e}{2e^2(m_i+m_e)n}J^2,
\quad\rho=(m_i+m_e)n.
\end{equation}
后一项是相对漂移能。若保留电子惯性却漏记相应能量与通量，可能把可逆漂移储能变化误归为加热。碰撞少时电子惯性提供交换和储存通道，本身不是正定熵源。

\section{电功怎样经过流能进入内能}
利用连续性将动量方程点乘本物种速度，再用
$-\vct u\cdot\nabla\cdot\vct P=-\nabla\cdot(\vct P\cdot\vct u)+\vct P:\nabla\vct u$，分别得到
\begin{align}
\partial_tK_i+\nabla\cdot(K_i\vct u_i+\vct P_i\cdot\vct u_i)
 &=en_i\vct u_i\cdot\E+\vct P_i:\nabla\vct u_i+\vct u_i\cdot\vct R_i,\\
\partial_tK_e+\nabla\cdot(K_e\vct u_e+\vct P_e\cdot\vct u_e)
 &=-en_e\vct u_e\cdot\E+\vct P_e:\nabla\vct u_e+\vct u_e\cdot\vct R_e.
\end{align}
因此压力功在流能和内能之间交换；电场功先进入物种总粒子能，而非直接全部进入温度。与Poynting定理相加后，电功也恰好抵消。封闭碰撞交换必须满足
\begin{equation}
\vct R_i+\vct R_e=0,\qquad
H_i+H_e+\vct u_i\cdot\vct R_i+\vct u_e\cdot\vct R_e=0.
\end{equation}
外部驱动、辐射、电离或开放边界另记，不能靠调整两物种热源使残差消失。

在准中性分解中$\E_H=\J\times\B/(en)$，所以$\J\cdot\E_H=0$。但$en\vct u_i\cdot\E_H$和$-en\vct u_e\cdot\E_H$可分别非零且互相抵消。这一恒等式只涉及欧姆律中特定分量，不允许把实测面内电场全部指定为Hall项。Hall场能重排物种能量路径，却不自动指定最终不可逆热化。

\section{七种量及边界通量}
\begin{longtable}{P{27mm}P{66mm}P{55mm}}
\toprule 名称 & 操作定义 & 易混淆之处\\\midrule\endfirsthead
\toprule 名称 & 操作定义 & 易混淆之处\\\midrule\endhead
物种电功 & $W_s=\int\dd t\int_Vq_sn_s\vct u_s\cdot\E\dd V$ & 含有序流储存和输运，不等于热。\\
流能 & $\int_VK_s\dd V$ & 需指定实验室系及物种平均流。\\
内能 & $\mathcal U_s=\int_V\operatorname{tr}\vct P_s/2\dd V$ & 含未分辨多束的二阶矩。\\
温升 & $\Delta T_s$或$\Delta T_{s,ab}$，指定位置和时间 & 密度和粒子交换改变温度与能量的对应。\\
焓流及压力输运 & 总粒子能流中的$U_s\vct u_s+\vct P_s\cdot\vct u_s$ & 标量时为$(U_s+p_s)\vct u_s$；各向异性时不能只用$5p_s\vct u_s/2$。\\
非热成分 & 指定核心拟合、阈值及参考分布后定义 & 五矩或十矩本身不唯一恢复尾部能量。\\
不可逆热化 & 指定碰撞、黏性、粗粒化或矩熵的定义及收敛 & 正压力功不保证正熵产，温升不保证热化。\\
\bottomrule\end{longtable}

由式\eqref{eq:Ui}积分得到
\begin{equation}
\Delta\mathcal U_i=\int\dd t\int_V(-\vct P_i:\nabla\vct u_i+H_i)\dd V
-\int\dd t\oint_{\partial V}(U_i\vct u_i+\vct q_i)\cdot\dd\vct A.
\end{equation}
内能账本中的对流能流和总能账本中的焓流是两种重排，不可混用后重复加$p\vct u$。$W_i/W_e$、$\Delta\mathcal U_i/\Delta\mathcal U_e$、$T_i/T_e$和所选不可逆热量比应分别输出；分母接近零或变号时应直接报告各项而非稳定“份额”。

\section{十矩新增的方程与热流闭合}
以下中心压力方程对离子、电子各有六个独立分量，分别取$(q_s,m_s,\vct u_s)=(e,m_i,\vct u_i)$和$(-e,m_e,\vct u_e)$：
\begin{align}
D_sP_{s,ab}+P_{s,ab}\partial_cu_{s,c}
+P_{s,ac}\partial_cu_{s,b}+P_{s,bc}\partial_cu_{s,a}
+\partial_cQ_{s,abc}\notag\\
=\frac{q_s}{m_s}\left(\epsilon_{acd}P_{s,cb}+\epsilon_{bcd}P_{s,ca}\right)B_d+C_{s,ab}.
\end{align}
磁项旋转压力张量但迹为零；取半迹返回内能式，其中$H_s=\operatorname{tr}\vct C_s/2$。新增的非对角压力参与欧姆律和应变功，新增的平行、垂直分量参与各向异性。但三阶$Q$仍未确定，十矩并非“无闭合”。

局地松弛常以$\partial_cQ_{s,abc}\simeq\nu_s(P_{s,ab}-p_s\delta_{ab})$建模，$\nu_s\sim v_{t,s}|k_{0s}|$。其迹为零，直接作用是各向同性化；即使矩熵增加，\textbf{该子步并不增加总内能}。热流在真实空间中搬运能量的作用不能仅由这样的零迹项代替。梯度闭合则可写成代表性形式
\begin{equation}
Q_{s,abc}=-\kappa_s n_s\,\mathcal S_{abc}(\partial_aT_{s,bc}),
\quad \kappa_s\sim\chi v_{t,s}/|k_{0s}|,\quad T_{s,ab}=P_{s,ab}/n_s.
\end{equation}
$\mathcal S$的平均或求和规范须按具体代码定义。常系数时$-\nabla\cdot Q$形成扩散，Fourier衰减随$k^2/|k_0|$，有别于局地松弛的固定衰减率。L256扩散的是压力偏离各向同性部分；Ng2020及后继温度梯度形式并不相同，不能只因都叫gradient closure而互换。

非局地Hammett--Perkins型算子近似$|k|v_t$相混响应，通常需Fourier或沿场非局地运算。适用的线性、磁化及几何排序不同；在强导场中用无磁化各向同性扩散可能过度跨场输运。闭合变得更复杂，并不自动保证压力正定、非线性稳定或真实碰撞熵产。\citep{Wang2015,Ng2017,L256,Ng2020,Closure2026}

\section{各向异性与多流体怎样改变加热解释}
旋转对称压力为$\vct P=p_\perp\vct I+(p_\parallel-p_\perp)\vct b\vct b$，内能$U=p_\perp+p_\parallel/2$，故
\begin{equation}
-\vct P:\nabla\vct u=-p_\perp\nabla\cdot\vct u-(p_\parallel-p_\perp)\vct b\vct b:\nabla\vct u.
\end{equation}
相同压缩率下，沿场应变使各向异性模型与标量模型有不同能量交换。无热流、磁化且满足双绝热排序时$D(p_\parallel B^2/n^3)=0$、$D(p_\perp/(nB))=0$，从而$T_\parallel\propto n^2/B^2$、$T_\perp\propto B$。这是可逆变形极限，不能把两个温度的分离直接命名为不可逆加热。俘获电子状态方程在此极限与通过粒子的等温响应之间连接，边界热库是必要信息。\citep{Wetherton2019,Wetherton2021}

再考虑同一物种分成群$a$，令$\vct w_a=\vct u_a-\vct u$。由展开中心矩得到
\begin{align}
\vct P&=\sum_a(\vct P_a+\rho_a\vct w_a\vct w_a),\\
Q_{ijk}&=\sum_a[Q_{a,ijk}+w_{a,i}P_{a,jk}+w_{a,j}P_{a,ik}+w_{a,k}P_{a,ij}
+\rho_aw_{a,i}w_{a,j}w_{a,k}].
\end{align}
即使每群内部热流很小，总三阶矩也可来自群间漂移。F635用这一思路给出电子多流体闭合；它增加的是电子群标签与相对流动，并非新增自洽离子加热。一个总电子流体在数学上仍可表示这些低阶矩，问题在于现有闭合是否能够从演化变量可靠预测它们。\citep{F635}

\chapter{问题、求解尝试与仍然缺少的环节}
\section{问题一：离子先热还是电子先热，为什么同一最终跳跃仍不够}
\subsection{已有尝试：共速双温慢激波}
Longcope与Bradshaw先取偏斜磁场重联后产生的流动，再以一维双温碰撞流体计算激波和传导前沿。此问题无电流的特殊约化使两物种共用密度和速度，电子惯性被忽略，电子压力梯度由电场平衡；它没有自行求解Hall扩散区。离子黏性把流能变成离子内能，库仑交换把能量送给电子，电子热传导又影响激波上游。\citep{L164}

共速、近各向同性时，可将两温方程示意写成
\begin{align}
\tfrac32nk_BD\theta_i+p_i\nabla\cdot\vct u
 &=\vct\tau_i:\nabla\vct u-\nabla\cdot\vct q_i+C(\theta_e-\theta_i),\\
\tfrac32nk_BD\theta_e+p_e\nabla\cdot\vct u
 &=-\nabla\cdot\vct q_e+C(\theta_i-\theta_e),\quad C>0.
\end{align}
这是展示各通道的简化式，不冒充原文所有系数。均匀无其他源时，温差$\delta\theta=\theta_i-\theta_e$满足$D\delta\theta=-4C\delta\theta/(3nk_B)$；因此温差松弛时间为$3nk_B/(4C)$。若文献以单一物种能量交换时间定义$C$，二者可能差因子2，比较时必须对齐定义。

原文计算分出离子黏性子激波、下游离子冷却区和电子传导前沿。早期离子温度可超出最终平衡温度，电子升温滞后；远下游平衡跳跃接近单温结果，却不代表中间路径相同。此结论超出H1的能力，但H2若使用同样两温、交换和输运方程，在适用极限内也应能够得到。因此这是对单温约束的修正，不是对全部Hall模型的否定。

\subsection{成立条件与缺口}
原文使用电子、离子Prandtl数约0.013、0.33，温度依赖的经典输运；稳态由射击法求解，另计算时变激波。部分强激波的早期离子热流接近或超过自由流量级，黏性修正甚至可能违反压力正定，恰好暴露经典离子闭合的边界。模型未包含X点直接电子加热，也未做实验两温的定量验证。不能用其电子热流较小的结果排除其他几何的直接电子加热。
\loc{L164}{本地45页版，式(4)--(7)、(14)--(18)；图4--7、9--14；物理页10、12、17、19--20、24、26、28、30、32、35、37。}

\section{问题二：双温Hall已有热点，完整双流体究竟增加什么}
\subsection{Browning与HIT-SI给出的基线}
Browning等的MAST磁环合并研究前列方程演化总压，后续图3才是初步独立两温扩展。离子黏性供热与电子电阻、超电阻供热产生不同空间图像，离子下游双峰与装置温升位置有联系。超电阻参数同时参与破冻结和电子热源，故热点并非仅由Hall项决定。本文未列该两温图专用网格和收敛序列，不可从总模型介绍补猜。\citep{L223}

HIT-SI的三维双温Hall模型更加完整地给出物种热源、交换和壁面热流。PSI-Tet与NIMROD中电子随电子流对流、压缩，欧姆热归入电子，黏性热归入离子；增强电子质量用于简化惯性项。模型能得到随注入频率变化的$T_i>T_e$，但装置功率包含注入器及反转过程，不能把全部热源归入单一X线。两代码比较还夹杂几何差别。\citep{Kaptanoglu2020}

\subsection{增加完整两流体后的可检验差别}
从H2到保留电子惯性的五矩，新增电子流能和快速电磁/电荷响应，但没有增加压力各向异性。若目标是远大于$\lambda_D,d_e$的下游温度，且两者采用一致热源和边界，这些新自由度可能只产生小修正；不能预设完整Maxwell一定改进加热。反之，在电子层，$D_e\vct u_e$和$\vct P_e$的非理想支撑可能改变电功位置。应通过局地欧姆项和能量项诊断确认哪一项造成差异。

Li与Livescu是很重要的反例：理论方程保留Braginskii物种输运，但实际四项验证强制$T_i=T_e$，删除独立离子能量演化。其六阶紧致差分、物理耗散和网格收敛说明数值框架的能力；并没有验证两物种加热分配。其所谓two-fluid又忽略电子惯性与位移电流，与通常“完整全Maxwell双流体”的定义不同。\citep{L247}
\loc{L223 / Kaptanoglu2020 / L247}{分别为图3（PDF7页）；式(1)--(13)、表I及图3--4（PDF3--4、8--9页）；第IV节关于$T_i=T_e$的实际测试约束与附录输运。}

\section{问题三：离子减速、压力支撑和温升是否由同一闭合解释}
\subsection{从局地松弛到非局地热流}
Wang等用五矩和十矩比较压力结构。十矩的有限压力松弛可以改善X点附近电子流及非对角压力，而强松弛近五矩会产生参照中没有的岛链。其图像常按相同重联磁通而非相同时间比较；$P_{xy}$峰值仍有显著误差。它证明新增张量自由度有用，并未证明温度、热流或全过程均正确。\citep{Wang2015}

Ng等2015年的Fadeev岛合并直接比较Hall和十矩。局地十矩改变离子压力支撑，却仍难得到参照中的大岛反弹与尺寸依赖。2017年非局地热流闭合使离子层厚从Hall的约$0.18d_i$改善到约$1.20d_i$，接近该算例参照的$1.15d_i$；但平均率的尺寸指数仍有偏差。\citep{Ng2015,Ng2017}

\subsection{温度梯度闭合的推进与未完成的加热账本}
Allmann-Rahn等2018年先检查热流散度，指出局地松弛的空间形状乃至符号可能错误，再以压力偏离各向同性的梯度闭合改善岛合并及反弹。其不同算例分别调节$k_0$，并非一套参数无校准成功。Ng等2020年改用温度梯度形式，平均重联率尺寸依赖和离子宽层进一步改善，并以同一典型$k_0d_s$测试重联与低混杂漂移不稳定性。\citep{L256,Ng2020}

这些研究对加热的意义在于：离子出流的减速长度、压力支撑和压力--应变交换位置会随闭合改变。\textbf{本报告的推论}是，仅用局地Hall温度峰去标定黏性，有可能掩盖错误的离子压力输运。原论文的主要成功指标却是层宽、反弹、率和不稳定性，而非经过实验验证的$Q_i^{\rm irr}/Q_e^{\rm irr}$。尚需把同一组计算中的完整压力功、热流、储能和边界项同时输出，才能闭合这一因果链。

成立条件包括二维岛合并几何、有限质量比和具体闭合波数。Ng2020指出电子各向异性结构仍不准确，跨场热流偏大；其附录还显示，对称稳态停滞点在零热流、无碰撞条件下可能趋向不物理压力。增加张量却继续设$Q=0$并不能保证稳态加热层可成立。
\loc{Wang2015 / Ng2017 / L256 / Ng2020}{Wang图4、8--10；Ng2017图1--5；L256图1--10与表1；Ng2020式(4)--(6)、图1--7、附录A。详细参数见附录。}

\section{问题四：电子各向异性有没有比标量温度更强的检验}
\subsection{可以进入流体方程的状态方程}
各向异性电子闭合将$p_{e\parallel},p_{e\perp}$写成局地$n,B$及外部来源参数的函数。这里的“局地函数”不意味着物理机制纯局地：外部热库和俘获/通过粒子的分界承担了非局地信息。新增的关键不是电子惯性，而是两个压力自由度和受俘获限制的热流响应。

Wetherton等2019年的MMS事件中，低密度区域电子接近等温Boltzmann响应，高密度区域更接近双绝热响应；一个固定标量电子压不能同时描述平行、垂直温度及转折。论文的数值对照使用H3D动理学离子加流体电子，必须归为混合模型；本报告仅将其用于验证可移植的电子闭合，不把离子结果算为双流体预测。\citep{Wetherton2019}

\subsection{从横向力平衡到有条件的温升标度}
该文强导场估计取峰值处$B\simeq B_g$，忽略电子压力对横向总压的主导作用，并取离子绝热$p_i\propto n^\Gamma$。由
\begin{equation}
p_{i\infty}\left[(n/n_\infty)^\Gamma-1\right]\simeq\frac{B_r^2}{2\mu_0}
\end{equation}
得$\widetilde n=(1+1/\beta_{ir})^{1/\Gamma}$，$\widetilde B=B_g/(B_r^2+B_g^2)^{1/2}$；$\beta_{ir}=2\mu_0n_\infty T_{i\infty}/B_r^2$。将二者代入电子状态方程，才得到标量温升。按原文低$\beta$渐近、$\Gamma=5/3$，
\begin{equation}
\frac{\Delta T_e}{m_iV_{Ar}^2}\simeq\frac{\pi}{36}\frac{T_{e\infty}}{T_{i\infty}}
\frac{B_r^2+B_g^2}{B_g^2}\beta_{ir}^{-1/5},
\quad T_{i\infty}\gg T_{e\infty},\quad B_g/B_r\gtrsim1.
\end{equation}
这里$m_iV_{Ar}^2=B_r^2/(\mu_0n_\infty)=2T_{i\infty}/\beta_{ir}$，因此指数由$-6/5$变为$-1/5$。系数$\pi/36$使用原文引用的渐近状态方程，本报告核验归一化及指数转换，不宣称重推其前序动理学积分。图2c的横轴是电子重联场beta，推导式中的$\beta_{ir}$是离子beta，不能混作同一参数。

原文典型范围得到约0.02的归一化电子温升；它是有条件的标度，不是电子占释放磁能的通用2\%，更不是不可逆熵产。电子温度的标度还借用了离子绝热闭合和峰值磁场假设，因此并非独立确定了离子加热。

2021年远下游研究进一步指出：局地磁俘获和有限通过时间造成的额外垂直增能，不能仅写成$n,B$的通用函数。它提出的能量修正需要轨道历史或额外状态信息；这正是有限局地闭合的可测边界。\citep{Wetherton2021}
\loc{Wetherton2019 / Wetherton2021}{前者PDF5--8页、图1--3与式(4)--(6)；后者PDF13--19页、表1、图6--10及式(4)。实验拟合条件见第4章。}

\section{问题五：出现三维不稳定性是否就解释了额外热化}
TenBarge等2019年以MMS 2015年10月16日事件参数计算全Maxwell五矩/十矩。十矩二维、三维的全域能量变化较近，离子内能增量较大；但电子各向异性峰虽接近观测幅度，峰值出现在错误的一侧。五矩还出现过强低混杂漂移活动与明显总能损失，不能把该模型的能量差直接解释为真实加热差。\citep{TenBarge2019MMS}

Allmann-Rahn等2021年及Bradshaw等2026年采用梯度型闭合，让十矩中的部分横向不稳定性得以增长。前者指出噪声幅度强烈影响启动与后续湍动；后者在类似MMS事件的三维非对称重联中得到局地闭合没有的漂移模、kink及通量绳。新增效果来自热流闭合对模态的阻尼选择，而非只来自“第三维”或“矩数更多”。\citep{E552,Closure2026}

Bradshaw等图6右还给出直接的反例约束：相对局地闭合，梯度闭合中的离子压力能和流能变化较小，电子压力能增长较快。恢复更多不稳定性并不自动带来更大的离子内能增量；这类全域能量曲线也没有分离不可逆熵产。

这两条线不能合并成一次独立重复：背景平衡、导场、质量比、闭合对象、数值扩散与种子均不完全相同。Bradshaw等的电子各向异性峰侧别仍不符，相关区域偏宽，且未给出完整三维网格/质量比收敛序列。\textbf{原文证明部分结构与模态出现，尚未证明其造成正确的两物种不可逆热化。}需要联合检验能量转移、热流、矩熵以及实验分布变化，而不能把“有湍流”直接当作耗散机制已知。
\loc{TenBarge2019MMS / E552 / Closure2026}{前者已读arXiv版图1及温度各向异性图；E552图4、6--9；Closure2026图1--11及附录A。正式版差异与收敛缺口保留在逐篇档案。}

\section{问题六：热流与多束有没有更合适的低成本表示}
\subsection{电子来源分群}
F635的两群电子十矩模型在固定离子背景上追踪电流片两侧来源。总热流由群内三阶矩、群间漂移和各群压力合成。它改善单群局地闭合的密度空洞与压力堆积，热流形状更接近参照，但幅度仍有偏差。这里“two-fluid”是两个电子群，\textbf{不是完整电子--离子双流体}；不能输出自洽离子加热。\citep{F635}

原文仍采用弱各向同性松弛；有限电子群数无法确保晚期每群保持单峰。图4在每个参照参数点取最佳时间相似度，不能读成两群流体在全参数扫描上的平均性能。图5色标不统一，也不能从图像近似断言幅度准确。其价值在于增加粒子来源历史这一低阶闭合缺失的信息；在本轮证据下，推广到离子参与、三维和实验仍属未验证。

\subsection{Landau/陀螺流体的时序检验}
L244从含电子两温、电子惯性和离子有限回旋半径的六场体系出发，但实际重联算例取冷离子子模型。Landau闭合降低平行温度扰动与增长率，使细尺度结构变平滑；按相同$\gamma t$比较的岛面积相似，不代表实际时间中的发展没有变慢。闭合还未表现出参照中的重联与加热延迟。\citep{L244}

这给出另一条超出H0/H1的途径：有限矩可近似部分平行相混，但自由能从解析变量消失，并不等于已在闭合中增加了相同数量的物种内能。若要称为加热预测，必须记录未解析自由能的去向或使用一致热库。实际冷离子算例没有回答离子、电子不可逆加热竞争；不能把理论六场变量清单当成已完成的双物种数值验证。

\subsection{为什么不按十三矩、二十矩的名称自动扩大结论}
独立演化三阶热流理论上能增加时间滞后和非局地响应，随后又需要四阶闭合、可实现性和稳定性控制。本轮检索发现相关会议摘要，但未取得足以核对离子--电子加热预算与实验的完整研究链，因此只列待核读，不将“矩数更多”写成已解决加热。Che的Buneman研究使用PIC诊断和近似流体分析，且算例未发展重联；它也不能替代独立流体重联验证。\citep{Che2014Buneman}

\chapter{能否与实验对上：逐装置审视观测链}
\section{匹配规则：先把模拟量变成测量量}
温度诊断不直接等于网格上的$\operatorname{tr}\vct P/(3n)$。多普勒线宽可能包含视线流速展宽、仪器展宽和杂质物种动力学；汤姆孙散射和航天器矩也有有限采样、背景与拟合响应。模拟至少应输出带空间、时间和视线响应的合成诊断。参数匹配还应包括$B_g/B_r$、$\beta_i,\beta_e$、质量比、$L/d_i$、密度非对称性、碰撞时间、驱动及壁面损失。

本文区分原文报告的误差与本报告建议的误差预算。原文没有给RMSE、协方差或独立验证集时，不根据图像补造百分比精度；图线重合也不等于参数无拟合的预测。

\section{MAST：热点位置与交换时间受到支持，热源仍不唯一}
Tanabe等测量强导场磁环合并。电子温度和密度来自300与130通道汤姆孙散射，空间分辨率约15与10毫米；离子温度来自32弦碳离子多普勒层析。二维剖面汇合多个放电及位置，不能当成单发全域瞬时测量。典型$B_t=0.3$--$0.8\,\mathrm T$、$B_r=0.07$--$0.15\,\mathrm T$，$d_i\sim0.1\,\mathrm m$，离子回旋半径小于约0.01米。\citep{Tanabe2015MAST}

\textbf{问题与尝试：}电子是否在重联中心局地升温，而离子在下游减速处升温？测量显示中心电子热点宽约2--5厘米，温度可在小于1毫秒内从约10升至200 eV；离子升温位于出流及下游。经过约4--10毫秒的物种能量交换时间，电子和离子都可出现三峰。这与双温流体的热源分离、交换和各向异性传导解释相容。

\textbf{成立条件：}Browning的图3是直接流体比较入口，但模型初始密度约$5\times10^{18}\,\mathrm{m^{-3}}$、温度10 eV以及耗散参数，并不构成对所有后续MAST放电的逐发匹配。其keV级离子峰与较低驱动实验的幅度有差别。实验导场增加使中心电子热点更尖，所测下游离子温度变化却较小，说明约束与热损失的影响不可遗漏。\citep{L223}

\textbf{剩余缺口：}缺少同一放电的完整电场、张量压力功、热流和磁能输入；原文还指出缺乏直接磁探针约束。因而位置符合支持“下游离子能量沉积/输运”这类解释，尚不能独立在黏性、压缩、多束展宽之间作唯一选择。其最强结论是空间、时序和输运约束，非普适物种份额。
\loc{Tanabe2015MAST}{PDF2--4页图1--5，核对几何、温度剖面、时间演化和导场扫描；L223的模拟温图在PDF7页图3。}

\section{HIT-SI：真实驱动、边界与两温耦合}
Kaptanoglu等把14.5--68.5 kHz注入频率、壁面条件及三维几何纳入双温Hall计算。离子Doppler与初步Thomson结果支持部分温度趋势，作者明确尚未进行严格定量验证。模型以离子黏性$550\,\mathrm{m^2/s}$、人工粒子扩散$250\,\mathrm{m^2/s}$及有效质量比100等为输入；壁面温度约3 eV，边界密度约$0.75\times10^{19}\,\mathrm{m^{-3}}$。\citep{Kaptanoglu2020}

\textbf{已解决到的程度：}双温模型能解释为什么独立$T_e$改变电阻率、进而反馈电流和驱动，同时离子可显著更热。图3把热源、交换和壁面流出放在同一预算内，比只展示$T_i/T_e$更有解释力。

\textbf{未解决：}两代码并非完全相同几何；人工粒子扩散尚未达到不敏感极限，中性、辐射和真实电路仍未完整纳入。注入器也在供热，因此积分功率不能直接命名为重联加热效率。已有勘误\href{https://doi.org/10.1063/5.0032632}{10.1063/5.0032632}的正文仍未取得，本报告不以未经更正核对的精确温度比例作为结论。这是实质版本缺口，不能用增加描述文字消除。

\section{MRX：能量账本能约束流体，但不是独立流体验证}
MRX使用磁探针重建磁场、电流，Mach探针约束离子流，Langmuir探针约束电子温度和密度，离子光谱诊断温度。电子流常由$\vct u_e=\vct u_i-\J/(en)$反演，电势和电场的估计也依赖测量与平衡关系。离子谱在有限视线与分辨率下仍须区分流动展宽；有关诊断原始论文和源数据未在本轮重分析。\citep{F679,F674,Yoo2017Heating}

\begin{longtable}{P{39mm}P{50mm}P{59mm}}
\toprule 指标 & 对称MRX（F674） & 非对称MRX（Yoo2017Heating）\\\midrule\endfirsthead
\toprule 指标 & 对称MRX（F674） & 非对称MRX（Yoo2017Heating）\\\midrule\endhead
离子电场沉积率 & $0.71\pm0.15$ MW & $0.44\pm0.09$ MW\\
电子电场沉积率 & $0.35\pm0.06$ MW & $0.35\pm0.07$ MW\\
报告中的热能变化项 & 离子$0.27\pm0.10$ MW，电子$0.26\pm0.10$ MW & 两者均约0.30 MW；不是由电功比直接给出\\
必要解释 & 热能变化项包含储存及边界焓流重排；更换控制体会变份额 & 密度非对称、分离面输运与控制体不同；不能外推对称2:1\\
\bottomrule\end{longtable}

F679还把有序离子加速、下游温升与有限碰撞联系起来，但其数值比较为PIC，不能归为双流体模拟成功。对流体研究最直接的要求是：同一控制体同时匹配物种电功、流能、温度和边界输运。只命中2:1，而没有命中近似1:1的热能变化项，实际上可能给出错误能量路径。

两项比值也不能假设测量误差独立后简单相除：它们共用密度、电场、控制体等测量输入，相关误差需要原始数据。原文部分热损失误差很大或只给上限。因此本报告保留各功率及其误差，不补造精确份额误差条。
\loc{F674 / Yoo2017Heating / F679}{F674正式版PDF6--9页、图6--8与表I；Yoo2017图10--11在PDF15页、式(10)在14页；F679图5、9--13及式(1)--(19)。}

\section{MMS近X线：电子闭合的可分辨预测}
\subsection{2019年观测与混合对照}
Wetherton等研究2015年10月31日约07:18 UT的MMS3穿越：距X点约$7$--$10d_i$，近对称且强导场，重联发生在相撞出流之间的磁通绳内部。电子分布矩与磁场沿轨迹给出$n,B,T_{e\parallel},T_{e\perp}$。作者采用$n_\infty=7.1\,\mathrm{cm^{-3}}$、$B_\infty=55.5$ nT、$T_{e\perp\infty}=32.4$ eV及$T_{e\parallel\infty}/T_{e\perp\infty}=1.22$描述来源。\citep{Wetherton2019}

\textbf{辨别力来自曲线转折：}图2a在密度约4--12 $\mathrm{cm^{-3}}$范围内，平行温度从近等温过渡为显著增加，而垂直温度只弱变；固定CGL或固定Boltzmann极限都不能覆盖整个区间。图1还比较俘获/通过边界和分布形状，因此并非只拟合单一温度峰。

混合模拟取$B_g/B_r=1.75$、总磁场归一电子$\beta_e=0.02$、$T_i/T_e=9.1$，力自由初态，$80\times40d_i$、$2048\times1024$网格、每格400个离子粒子，以显式电阻破冻结。在$t=80/\Omega_{ci}$选取宽约$3.5d_i$的出流截面与观测比较。低质量比100的PIC在密度空洞处出现额外平行温度下降；电子流体闭合没有这一明显偏差。这里是有限电子通过时间排序的比较，不能推成所有流体优于真实质量比动理学。

\textbf{验证强度：}原文称定量符合，但没有在本文提供完整残差统计、误差协方差或混合模型网格收敛序列；截面按观测宽度选择。它比仅观察热点位置强，但仍是单事件、带匹配条件的闭合支持。离子是动理学粒子，不属于完整双流体分配检验。FPI电子分布和磁场测量提供数据基础，本文未详列全部仪器响应和校准误差，本报告不凭常规仪器指标补充成本文实际分辨率。

\subsection{与2015年10月16日十矩事件不可混为同一次验证}
TenBarge2019MMS和Closure2026使用的是另一个非对称、弱导场事件。局地十矩的电子各向异性峰幅约2.5虽近观测量级，位置在磁鞘侧，观测却在磁层侧；后继梯度闭合仍有侧别和过宽区域问题。两事件的导场、非对称性和电子磁化条件不同，不能从强导场EoS的成功直接宣称弱导场十矩通过实验。\citep{TenBarge2019MMS,Closure2026}

\section{MMS远下游：参数拟合、外部压缩和闭合失败同时存在}
Wetherton等2021年分析2016年1月21日约01:06 UT的MMS2事件，约$100d_i$下游、$B_g/B_r\sim0.7$、入流略不对称。这个距离来自前序时序分析，不能当作局地直接测距。作者为两侧入流和出流分别选择外部来源参数，使状态方程匹配温度曲线；参数见表。\citep{Wetherton2021}

\begin{center}
\begin{tabular}{lrrrr}\toprule
区域 &$n_\infty$（cm$^{-3}$）&$B_\infty$（nT）&$T_{\perp\infty}$（eV）&$T_{\parallel\infty}/T_{\perp\infty}$\\\midrule
入流1 &47&30&28&0.9\\
出流 &41&29&30&1.0\\
入流2 &40&41&30&0.9\\\bottomrule
\end{tabular}
\end{center}

这些是原文表1的模型来源参数，\textbf{不能写成三个远处热库都被独立直接测量}。图6多数区间温度随局地$n,B$变化被解释，图7有三条区域特征曲线；但若干紫色区间垂直温度显著超过原闭合预测。

作者把偏差联系到局地俘获、有限通过时间与整个区域的磁压缩，增添$\Delta\mathcal E=\mu\int\partial_tB[\vct x(t)]\dd t$。深俘获例的推断能量增强约1.35，与选定$B/B_\infty\simeq1.34$相符；同时出流区额外项通常较小，最强垂直增能位于部分入流。图9--10的增能是依据分布模型反演，不能当成直接测得的独立能量通量。

\textbf{本轮最重要的限制证据：}原文明确没有$\Delta\mathcal E$对$n,B$的形式依赖，因而不能把这一步视为已完成的通用改进状态方程；压缩起源及是否延伸至X线也未确定。已有曲线符合支持“来源与俘获调节温度”，但经调参后的符合、反演增能和无参数预测是三种不同证据。垂直温升也可能来自区域压缩，不能全部记在重联局地耗散账下。

\section{实验结论矩阵}
\begin{longtable}{P{23mm}P{46mm}P{40mm}P{33mm}}
\toprule 证据 & 已支持的量 & 未通过的强断言 & 当前程度\\\midrule\endfirsthead
\toprule 证据 & 已支持的量 & 未通过的强断言 & 当前程度\\\midrule\endhead
MAST & 温度热点、交换时序、导场趋势 & 黏性热源唯一；通用分配 & 定性空间/趋势\\
HIT-SI & 两温趋势与驱动耦合 & 严格定量温度比；仅X点供热 & 条件性定性\\
MRX & 分物种电功与有限体积预算 & 2:1不可逆热化比；纯流体已验证 & 实测预算有误差\\
MMS强导场 & 各向异性温度及密度转折 & 完整双流体两物种分配 & 事件级闭合支持\\
MMS远下游 & 分区域EoS曲线和偏离位置 & 来源无拟合；局地通用垂直热源 & 拟合约束与失败并存\\
MMS弱导场十矩 & 部分结构、峰幅及能量演化 & 正确热点侧别；热化比获证 & 存在实质不吻合\\
\bottomrule\end{longtable}

\chapter{对Hall结果的判决及下一步研究}
\section{模型能力比较：新增自由度、实际成果与代价}
\small
\begin{longtable}{P{26mm}P{39mm}P{43mm}P{34mm}}
\toprule 模型 & 新增可计算量 & 已有直接证据 & 仍受限制\\\midrule\endfirsthead
\toprule 模型 & 新增可计算量 & 已有直接证据 & 仍受限制\\\midrule\endhead
共速双温 & 两个温度、交换滞后、传导前沿 & L164的碰撞激波时空结构 & 给定出流，不求扩散区；经典离子闭合可失效\\
双温耗散Hall & 独立物种预算、壁面热损与电子对流 & L223、HIT-SI的温图和趋势 & 热源分配/输运系数/边界；非完整张量\\
惯性五矩 & 电子流能、惯性响应、两套标量总能 & Wang2015及MMS五矩对照 & 无各向异性；MMS例能量损失较大\\
局地十矩 & 六分量压力、非对角支撑、各向异性 & Wang2015；TenBarge2019MMS全域热能 & 零迹松弛难替代真实热输运；峰侧别错\\
非局地十矩 & 近似尺度依赖相混及压力输运 & Ng2017离子层与压力支撑改善 & 无磁化闭合边界；电子结构仍偏差\\
梯度十矩 & 热流散度、岛反弹、部分横向模态 & L256、Ng2020、E552、Closure2026 & 参数校准、跨场扩散、正定与三维收敛\\
各向异性EoS & 给定来源下平行/垂直电子温度 & Wetherton2019/2021观测检验 & 未提供通用离子闭合；局地俘获历史遗漏\\
Landau/陀螺流体 & 部分平行相混、惯性与有限回旋半径 & L244的增长率/细结构变化 & 实际冷离子；热化延迟与自由能归宿\\
来源分群电子 & 群间漂移对总压力和热流的贡献 & F635电子独有重联示例 & 固定离子；群标签、晚期多束及幅度偏差\\
十三/二十矩 & 理论上演化热流并闭合四阶矩 & 本轮只追到待核查线索 & 不据矩数或会议摘要声称加热验证\\
\bottomrule\end{longtable}
\normalsize

\section{三个比较问题的逐项判决}
\subsection{能得到Hall得不到的结论吗}
\textbf{相对H0/H1：能。}独立能量使温度比例和滞后不再由输入规定；L164与HIT-SI实际展示了重要差别。新增方程是两套热力学，而不是Hall电磁项。

\textbf{相对标量H2/H3：部分能。}十矩能描述非对角压力、各向异性温度和相应应变功；EoS能预测由俘获限制热流所造成的温度--密度转折。这些是标量压力体系无法直接输出的量。其正确性仍取决于闭合，尤其是热流。模型能表示某一量与模型已经准确预测该量必须分开。

\textbf{相对H4：没有由名称保证的优势。}如果Hall约化已保留相同压力张量、热流和独立能量，区别主要转到电子惯性、电荷分离、位移电流及高频尺度。在低频准中性适用区，完整Maxwell多流体应逼近相应Hall极限。大幅偏差首先要求核查参数、边界、能量一致性与数值误差，而非自动视为新物理。

\subsection{能和实验对上吗}
可以对上部分温度形态、趋势和电子闭合关系；本轮证据没有支持“完整双流体已同时精确匹配实验$T_i,T_e$、张量压力、热流、两物种不可逆加热和总能预算”。实验验证最强的部分集中在有限控制体预算或特定电子状态方程，而不是一套跨装置闭合。原文未给统一误差指标，不能以本报告的表格整齐代替这一缺口。

\subsection{能独立印证Hall加热结果吗}
支持程度必须按独立性排序：同一模型的两份代码可检查实现；不同模型共用同一热源分配和调参，只能提供有限的稳健性证据；真正不同闭合在未用于校准的观测量上同时成功，才对物理解释有更强支持。L164与MAST下游离子先热的路径相容，但碰撞性、几何和供能方式不同，不能当作同参数重复验证。十矩闭合的离子层变化与MMS各向异性偏差则限制把Hall峰温相符等同于正确机制的做法。

\section{已解决与未解决：四类缺口分别处理}
\begin{longtable}{P{40mm}P{56mm}P{52mm}}
\toprule 问题 & 已有求解尝试及条件 & 剩余缺口与类别\\\midrule\endfirsthead
\toprule 问题 & 已有求解尝试及条件 & 剩余缺口与类别\\\midrule\endhead
温差怎样产生与消失 & 两温能量、碰撞交换、热传导；L164/HIT-SI & 快交换与无碰撞过渡；输运闭合可信度\\
离子出流怎样变成压力 & 十矩离子压力支撑、岛合并反弹；Ng/L256 & 缺实验压力功和热流联合验证；证据不足\\
电子为何平行更热 & 俘获EoS，强导场MMS & 弱导场非旋转对称压力与局地历史；模型边界\\
三维波动是否改变分配 & E552/Closure2026恢复部分模态 & 未给收敛物种熵产；目标尚未检验\\
多束与热核如何分开 & F635按电子来源分群 & 离子、三维及动态分群未验证；表示限制\\
Hall能否得到通用加热比 & 装置双温Hall与实验预算 & 源、边界、几何依赖且观测定义不同；无此证据\\
更多矩是否更好 & 原理上可演化热流 & 十三/二十矩完整证据未取得；文献覆盖缺口\\
原始各向异性流体链 & Ohia2012/2015原始论文已定位 & 正文/图式访问缺口未闭合，不据摘要入结论\\
\bottomrule\end{longtable}

前六项是本报告依据跨文献作出的综合判断，部分与原文明确限制重合；最后两项是本轮证据覆盖的缺口。未取得原文不能写成学界没有做过，模型不能表示尾部也不意味着所有加热问题都需要全动理学。

\section{研究设计一：先把模型差异归因到一条新增方程}
选择两个可控问题：无导场Fadeev岛合并用于离子压力支撑；强导场近对称重联用于电子温度各向异性。第一阶段不做装置全尺度，优先复用已公开输入，逐项核对归一化、边界和扰动，而非只抄论文表格。

同一初态与同一边界下依次计算：H1单热力学、H2独立能量、H3电子惯性、五矩全Maxwell、局地十矩、梯度十矩；强导场支线再与EoS比较。H1需明确指定分配比例，不能把该比例与其他模型的输出混合统计。H3与全Maxwell的差异要单独隔离；若同时换代码，应先做共享物理极限验证。

\textbf{消融顺序：}固定电子惯性，改变标量/张量压力；固定压力模型，改变热流；固定闭合参数，改变质量比和光速；最后才改变碰撞、导场和几何。碰撞交换、黏性及电阻必须在能量中成对记账。对局地松弛，额外输出它对压力迹的零贡献与对矩熵的贡献，避免把两者混同。

\textbf{输出：}各物种$W_s,\Delta K_s,\Delta U_s$、压缩功、无迹应变功、边界焓流、热流、碰撞交换、物理耗散、数值残差；$T_{\parallel},T_\perp$、压力特征值、热点位置与宽度。分别按相同时间和相同重联磁通比较，以区分演化时序与状态结构。

\textbf{成败门槛（本报告建议，非文献已有结果）：}三档网格及至少两档时间步，在最细两档上主要积分预算和热点宽度变化小于5\%；全能量残差与分物种预算残差各小于同期释放磁能的1\%，并小于所讨论模型差值的五分之一；压力保持正定，限制器/下限带来的能量改变量单列。若差异小于数值不确定度，不报告“新增物理改变加热”。这些门槛应在试算前固定，必要时因实验精度调整并留记录。

\section{研究设计二：不用同一条温度曲线同时调参与验证}
选择MAST的一组可获取放电或MMS一个公开事件。MAST需在模拟中实现实验驱动和热损边界，合成碳离子谱与电子散射响应；MMS需沿候选轨迹采样$n,B,T_{e\parallel},T_{e\perp}$，保留航天器时间到空间映射的不确定性。原始数据未取得前不得承诺已可重现。

先划分校准与验证对象：只用上游$n,B,T_s$及几何选择边界；若必须用一个温度剖面拟合$k_0$、黏性或来源参数，另一个时刻、位置、导场工况或热流曲线必须留出。固定参数后再预测测试集，不允许逐事件继续调参后声称通用。

对Wetherton2021类三来源拟合，先用入流及独立磁场资料限制来源，盲预测出流温度及紫色偏差区。若只有逐区域自由选择四个参数才能吻合，应报告拟合解释能力；若加入来源记忆变量后能预测未参与拟合的垂直增温，才支持模型扩展的预测价值。

\textbf{成败判据：}模拟与观测在预先指定区间的差值是否落在测量、参数传播与数值误差合成区间内；热点侧别和温度转折是否正确；能量残差是否明显小于待解释热量。原文缺少可用协方差时先报告分项误差，不机械计算约化$\chi^2$。若温度命中而功率/通量失配，判为机制尚未通过。

\section{研究设计三：区分净增温、散射与不可逆性}
对十矩模型，在总温度匹配后，进一步检查压力特征值、各向异性、热流与熵闭合。对正定压力张量，局地Gaussian矩熵每粒子可写为
\begin{equation}
\sigma_G/k_B=\frac12\ln\!\left(\frac{\det\vct P}{n^5}\right)+\text{含单位基准的常数}.
\end{equation}
这是用有限矩重建的Gaussian熵，不是未知真实分布的精细熵。各向同性化在固定迹下提高行列式，可提高$\sigma_G$却不增加内能；此诊断能揭示“散射增加”与“加热增加”的混淆，但不能替代分布测量。

对来源分群模型，同时报告群内内能和群间漂移贡献，检查总矩合成前后的能量恒等式。若温升主要来自群间漂移，论文和报告都应明确为有效二阶矩升高；需要实验可分辨的多束或碰撞松弛证据，才能进一步论证热化。此阶段只在关键点使用少量动理学参照，不扩展成本轮研究主线。

\section{最小可复查交付清单}
每个算例保存初边值、单位、质量比、$c/V_A$、网格、步长、离散格式、闭合参数及限制器，列出物理参数与校准参数。运行记录须保留能量残差、压力最小特征值、源项积分及边界通量。最终结论按“新增方程改变了哪个可测量、幅度超过何种误差、在哪些参数内成立”撰写；失败工况与成功工况一起保留。此报告完成的是文献和研究设计，没有声称执行以上模拟。

\appendix
\chapter{逐篇证据与数值实现档案}

本附录将本专题采用的证据按问题次序列出。旧记录来自已有全文笔记，本轮定向复用不重复计数。新增两篇为电子闭合的观测检验，其中数值部分含混合模型；不能称为新增两次完整双流体模拟。

本专题采用21篇核读证据，其中19篇复用既有记录、2篇新增核读；新增主版本共43个PDF物理页（含封面、参考文献及重复图页），全库144增至146篇。页面数量不等于独立研究数量。


\section{L164：2010年}\label{paper:L164}

\textbf{Slow shocks and conduction fronts from Petschek reconnection of skewed magnetic fields: two-fluid effects}\citep{L164}

\textsf{本轮状态：}已有全文记录复用，不增加计数。

\textsf{核读范围：}PDF 1–45 including all Sections 1–7 and references;  no internal appendix\par\textsf{图式：}Direct PDF visual: Eqs.4–7 (p10),  Eqs.14–18 (p12),  Figs.4–7, 9–14 (pp17, 19, 20, 24, 26, 28, 30, 32, 35, 37). Remaining figures/captions read in text.

\href{run:../../papers/L164.pdf}{本地原文PDF}\quad\href{https://doi.org/10.1088/0004-637X/718/2/1491}{DOI/版本入口}

\paragraph{实际模型与问题}采用准中性、共同密度和共同速度的一维两温碰撞流体描述平行激波，再映射到外部给定的Petschek偏斜磁场重联出流。共同速度来自本问题无电流的特殊约化，并非一般Hall假设。保留离子黏性、两物种经典热导及库仑温度交换；电子惯性忽略，压力梯度电场维持电子力平衡。它不是自洽Hall扩散区模拟。

\paragraph{参数、初边值与实现}例子包含入流马赫数2.4、3.6、5，初始双曲正切速度层宽1–500个归一化离子黏性长度；电子和离子Prandtl数0.013、0.33，交换参数0.019。输运系数随温度的5/2次方变化，交换频率随nTe\textasciicircum{}(-3/2)变化。中心镜面、远端持续入流；稳态用射击积分，时变网格和时间步未报告。偏斜角约≤120°时几何约化较可靠。

\paragraph{求解得到什么}离子黏性子激波先把出流动能转为离子内能，随后长下游区域的碰撞交换加热电子，电子传导又形成前导热波；三个区域及相关时间尺度明显分离。早期离子温度可超过最终平衡值，电子升温滞后，而远下游最终跳跃仍可与单温结果一致。初始速度层足够宽时，局地两温大小关系也可改变，不能压成固定加热比。

\paragraph{成立条件与剩余缺口}原文发现部分强激波早期离子热流达到或超过自由流值，黏性修正甚至违反压力正定条件，故经典闭合恰在关键加热层可能失效。电子热流较小与未包含X点直接电子加热有关，不能外推到所有重联。模型不含自洽Hall场、完整动理学尾、动态磁通管截面和反平行极限，也没有进行实验Ti/Te定量验证。

\paragraph{对Hall比较的意义}这是两温流体可研究热化过程的清楚实例：它把离子流能耗散、物种交换和电子传导在空间与时间上分开，同时主动揭示经典离子闭合的失效。可据此说明“双流体并非不能研究加热”，但不能把给定出流的碰撞激波结果当作无碰撞Hall重联的普适分配律。文中与拾取加热等动理学机制的联系应作为模型边界保留。

\section{L223：2016年}\label{paper:L223}

\textbf{Two-fluid and magnetohydrodynamic modelling of magnetic reconnection in the MAST spherical tokamak and the solar corona}\citep{L223}

\textsf{本轮状态：}已有全文记录复用，不增加计数。

\textsf{核读范围：}PDF 1–16: all body,  methods,  results,  discussion,  figures/captions,  references;  no internal appendix\par\textsf{图式：}Direct PDF visual: Eq.1–5 and definitions (pp3–4),  Fig.3 (p7),  relaxation equations/Figs.4–8 (pp9–10, 13–14);  remaining figure captions read in full text

\href{run:../../papers/L223.pdf}{本地原文PDF}\quad\href{https://doi.org/10.1088/0741-3335/58/1/014041}{DOI/版本入口}

\paragraph{实际模型与问题}采用二维直角和紧环径比轴对称环面Hall MHD描述同向电流磁通绳合并，保留电子漂移、电子标量压力、电阻、超电阻及热传导。前列方程仍演化总压，图3才给独立离子、电子温度的初步扩展：离子由黏性供热，电子由电阻和超电阻供热。后半篇的Taylor弛豫计算只估计磁能释放，不能混作物种加热模拟。

\paragraph{参数、初边值与实现}基准密度5×10\textsuperscript{18} m\textsuperscript{-3}、长度1 m、磁场0.5 T，初始Ti=Te=10 eV，归一化离子惯性长度0.145。典型电阻10\textsuperscript{-5}、离子黏性10\textsuperscript{-3}、平行电子热导0.1、垂直离子热导10\textsuperscript{-7}；图3超电阻10\textsuperscript{-8}。使用HiFi及前序模型，但本文没有报告该两温加热图专用的网格、时间步或收敛检查，完整两温方程亦未列出。

\paragraph{求解得到什么}图3的初步两温环面计算产生下游离子双热峰，峰值可达keV量级，与早期MAST高驱动观测有定性联系；电子热源主要跟随超电阻涉及的电流梯度，分布不同于离子。中心或空心电子温度与岛排出的联系仍是初步解释。弛豫模型另得到几何相关磁能释放及近似极向磁场平方标度，但不能给出热量物种比例或空间分布。

\paragraph{成立条件与剩余缺口}超电阻是自由模型参数，并直接影响电子加热位置；有限标量闭合不描述完整粒子分布或非热尾。离子热峰高于后来较低驱动实验，实验比较仍属初步，不能称为加热比例定量验证。Stanier2014学位论文、源代码和独立数据未核读，本文也未给两温热图的专门收敛证据；没有推导普适Qi/Qe。

\paragraph{对Hall比较的意义}本篇直接证明两温Hall扩展能够在指定耗散、传导和几何条件下计算离子与电子局地热图。证据应表述为“模型能做条件性热预算和实验比较”，不能表述为“Hall项本身提供不可逆热化”或“物种分配已经普遍解决”。同时必须区分前列总压方程、后续初步两温计算和独立磁能弛豫估计三种证据。

\section{Kaptanoglu2020：2020年}\label{paper:Kaptanoglu2020}

\textbf{Two-temperature effects in Hall-MHD simulations of the HIT-SI experiment}\citep{Kaptanoglu2020}

\textsf{本轮状态：}已有全文记录复用，不增加计数。

\textsf{核读范围：}PDF 1–35: all body,  methods,  results,  conclusions,  references and appended repeated/enlarged figure pages;  no internal appendix.\par\textsf{图式：}Direct PDF visual: Eq.1–13/Table I (pp3–4),  Figs.2–4 (pp7–9),  Figs.5–13/scan results (pp11–17);  overview of appended repeated images pp19–35. Key Ti/Te map and heating power plot inspected. Some noncentral scan figures reviewed at contact-sheet scale.

\href{run:../../papers/Kaptanoglu2020_accepted.pdf}{本地原文PDF}\quad\href{https://doi.org/10.1063/5.0006311}{DOI/版本入口}

\paragraph{实际模型与问题}在PSI-Tet和NIMROD中采用三维两温Hall MHD，独立演化Ti、Te，欧姆定律保留Hall漂移及增强电子质量的简化电流时间导数惯性项。电子温度随电子流对流和压缩；电子得欧姆热、离子得黏性热，两者通过库仑交换及各向异性热传导联系。闭合仍是标量压力、Spitzer电阻和常数各向同性黏性，另有人工粒子扩散。

\paragraph{参数、初边值与实现}注入频率14.5–68.5 kHz，边界密度0.75×10\textsuperscript{19} m\textsuperscript{-3}、温度3 eV，黏性550 m\textsuperscript{2}/s、扩散250 m\textsuperscript{2}/s，有效质量比100。PSI-Tet用103887四面体、2.8 cm三阶单元，等效9 mm，Crank–Nicolson步长≤40 ns；NIMROD用1.8 cm四阶极向单元和11环向模。两代码注入器几何不同，另扫描壁温、密度、扩散和增强质量。

\paragraph{求解得到什么}独立两温改变了单温模型的温度、密度与电流演化；较高驱动频率下离子明显热于电子。图4显示两物种相似的大尺度空间结构，但温度幅度不同；图3列各物种热源、交换和壁面热流，积分范围包括注入器。作者称与离子Doppler和初步Thomson数据定性一致，但明确未做严格定量验证。压缩功也不能全部视为不可逆热量。

\paragraph{成立条件与剩余缺口}温度依赖壁面条件和人工扩散，尚未达到对扩散系数不敏感的收敛极限；真实驱动电路、中性/原子/辐射过程及更真实黏性仍缺失。已有勘误DOI10.1063/5.0032632经元数据确认，但更正全文未取得；本次只保证35页接收稿完整核读，不宣称掌握更正后的全部公式或精确比值。两个代码的差异也不全来自两温处理。

\paragraph{对Hall比较的意义}这是本分支取得的最直接装置尺度两温Hall实例，证明可同时计算独立温度、局地热图和带输运损失的物种预算。但结果受驱动、边界与耗散闭合约束，且部分功率来自注入器反转过程，并非仅重联X点。应采用定性和条件性结论；未核勘误、人工扩散敏感性及缺失电路模型使其不能支持参数无关的Qi/Qe或精密实验验证。

\section{L247：2019年}\label{paper:L247}

\textbf{High-order two-fluid plasma solver for direct numerical simulations of plasma flows with full transport phenomena}\citep{L247}

\textsf{本轮状态：}已有全文记录复用，不增加计数。

\textsf{核读范围：}PDF1–44: entire text and references including Appendix A plasma transport,  B quasineutrality,  C single-fluid limits.\par\textsf{图式：}Direct visual Eq15–19 (p11),  identical-temperature restriction (p16),  Figs1, 3, 4, 7, 9, 10, 11 (pp18, 21, 25, 27, 35),  AppendixA EqsA1–A9, A15–A19 (pp31, 33). All other figures/captions read in text.

\href{run:../../papers/L247.pdf}{本地原文PDF}\quad\href{https://doi.org/10.1063/1.5082190}{DOI/版本入口}

\paragraph{实际模型与问题}构建准中性、忽略电子惯性与位移电流的Braginskii两流体方程，保留独立物种内能、磁场方向相关黏性和热流、摩擦、热力、电阻加热及温度交换；附录列出输运系数和单流体极限。但第IV节明确四项实际演示均强制Ti=Te，并删去独立离子能量演化，实际测试的是单温扩展Hall系统，不能当作两温分配验证。

\paragraph{参数、初边值与实现}波测试用384点；激波Rei为6.3×10\textsuperscript{3}和2.3×10\textsuperscript{4}。重联di/L=1、Rei约124–850、Rem约58–890，主例426/112；参考长度10 m、密度10\textsuperscript{15} m\textsuperscript{-3}。CFDNS用六阶紧致差分和物理输运；主例约256×512收敛，以768×1536为参考，观测阶数6.58。正文未给具体时间步，另有更高Re和激波网格扫描。

\paragraph{求解得到什么}在所选物理输运系数下，波色散、激波和重联基准得到高阶网格收敛；局地共同温度图显示磁能转化及热扩散的重要性。激波误差阶数约6.08–6.14，重联磁通在所试较大黏性和磁雷诺数处趋于弱敏感。论文提供了完整物种热方程及物理正则化的实现路径，但没有给出独立Ti/Te、交换过程或加热比例的验证结果。

\paragraph{成立条件与剩余缺口}最重要边界是所有验证算例固定Ti=Te，不能从完整方程形式推断已演示双温预测。经典碰撞输运在冠层扩散层可能不满足局地平衡排序；磁通及饱和时间随网格或Re收敛，也不能证明微观热化和物种分配正确。零电子惯性等约化排除相关电子尺度；原始输运推导、CFDNS源代码及独立双温测试未复核，也无实验温度比对照。

\paragraph{对Hall比较的意义}应将本篇用作能量方程、各向异性输运和数值收敛的方法证据，尤其支持区分物理输运与隐含数值耗散。作者关于高Re宏观结果可代表无碰撞重联的讨论，只由其实际磁通等指标检验，不能扩张为无碰撞加热闭合获验证。全文44页含附录均已读；正式出版年为2019，2018是预印本日期，两者不是独立文献。

\section{Wang2015：2015年}\label{paper:Wang2015}

\textbf{Comparison of multi-fluid moment models with particle-in-cell simulations of collisionless magnetic reconnection}\citep{Wang2015}

\textsf{本轮状态：}已有全文记录复用，不增加计数。

\textsf{核读范围：}1-15（arXiv v3全部）\par\textsf{图式：}图1-10及式20-24、36；图4、8-10视觉核验

\href{run:../../papers/Wang2015.pdf}{本地原文PDF}\quad\href{https://doi.org/10.1063/1.4906063}{DOI/版本入口}

\paragraph{实际模型与问题}五矩和十矩全Maxwell，局地Hammett-Perkins松弛；k0e=k0i=1/de为targeted，10\textasciicircum{}4/de近五矩。

\paragraph{参数、初边值与实现}主Harris例Lx×Ly=100×50di, a=.913di, nb=.3, Ti/Te=5, mi/me=25, vA/c=1/15, 扰动.1；流体4095×2047（约8格/de），PIC5760×2880（约12格/de），每物种2.7e9粒子。另有mi/me=25, 100, 256, 400质量比扫描。

\paragraph{求解得到什么}图4通量时间史明显不同；同重联磁通下targeted十矩的X点附近电子流和非对角压强较接近PIC，五矩及强松弛出现PIC没有的磁岛链。质量扫描给出电子层长约(me/mi)\textasciicircum{}(3/8)、厚约(me/mi)\textasciicircum{}(1/3)。

\paragraph{成立条件与剩余缺口}对比结构常取不同时间（图10 PIC75、流体132.5）；Pxy峰值PIC约.024对流体.006，标量化仍偏强；下游欧姆项峰位不一致、抵消不足。真实质量比层长12-17de和厚.4-.6de是外推，不是直接计算；无完整热流张量检验。

\paragraph{对Hall比较的意义}能复现部分局地压力支撑，不能写成重联全过程或所有压力分量定量吻合；局地闭合参数仍需选择。

\section{Ng2015：2015年}\label{paper:Ng2015}

\textbf{The island coalescence problem: Scaling of reconnection in extended fluid models including higher-order moments}\citep{Ng2015}

\textsf{本轮状态：}已有全文记录复用，不增加计数。

\textsf{核读范围：}1-8（全部）\par\textsf{图式：}图1-6文字及图注；主要结果图已视觉复核

\href{run:../../papers/Ng2015.pdf}{本地原文PDF}\quad\href{https://doi.org/10.1063/1.4935302}{DOI/版本入口}

\paragraph{实际模型与问题}电阻MHD、Hall MHD电阻／超电阻、十矩局地松弛对照PIC；Fadeev岛合并。

\paragraph{参数、初边值与实现}epsilon=.4, nb=.2n0, Lx×Lz=4piλ×2piλ, x周期, z导体／反射；mi/me=25，Ti=Te，beta=1；平均窗0-1.5tA，模拟至2.5tA。Hall eta=.009或etaH=1e-4及eta=1e-5；mu=.01。

\paragraph{求解得到什么}Hall区峰值／平均率对系统尺寸仅弱敏感，约λ\textasciicircum{}-0.2至\textasciicircum{}-0.25；十矩调低ki到1/(3di)改善小系统离子压力支撑，但大系统次级岛不断抛出，合并完成且缺少PIC弹回。

\paragraph{成立条件与剩余缺口}本文比较PIC平均率约λ\textasciicircum{}-0.8；引言早期研究的\textasciicircum{}-0.5不可混用，盒子宽高比不同。十矩ki是可调参数，出现PIC中没有的次级岛；本文不是十矩全面成功的证据。

\paragraph{对Hall比较的意义}必须比较同几何、同平均窗；Hall率可高2-3倍。

\section{Ng2017：2017年}\label{paper:Ng2017}

\textbf{Simulations of anti-parallel reconnection using a nonlocal heat flux closure}\citep{Ng2017}

\textsf{本轮状态：}已有全文记录复用，不增加计数。

\textsf{核读范围：}1-5（全部）\par\textsf{图式：}图1-5文字及图注；主要结果图已视觉复核

\href{run:../../papers/Ng2017.pdf}{本地原文PDF}\quad\href{https://doi.org/10.1063/1.4993195}{DOI/版本入口}

\paragraph{实际模型与问题}全Maxwell十矩；非磁化三维Hammett-Perkins热流Fourier闭合，chi=sqrt(4/(9pi))；有限体积输运、局地隐式源与FFTW。

\paragraph{参数、初边值与实现}Fadeev岛、epsilon=.4, nb=.2, Lx×Lz=4piλ×2piλ, mi/me=25, vt, e/c=.35, Ti=Te, beta=1, ωpe/Ωce=2，10\%磁扰动；平均0-1.5tA。

\paragraph{求解得到什么}非局地闭合得到1.20di片宽，PIC1.15di，Hall .18di；离子压力支撑区与速度、非旋转对称度改善；平均率λ\textasciicircum{}-0.45比局地\textasciicircum{}-0.2更接近PIC\textasciicircum{}-0.8。

\paragraph{成立条件与剩余缺口}仍未复现PIC尺寸标度，离子非旋转对称度偏低；次级岛影响速率；电子各向异性偏弱；非磁化闭合不应直接外推强导引场。Langmuir阻尼测得.044ωpe，精确Vlasov .034ωpe，存在约29\%误差。

\paragraph{对Hall比较的意义}成功是特定离子结构改善与减少调参，不是所有动力学完全重现。

\section{L256：2018年}\label{paper:L256}

\textbf{Temperature gradient driven heat flux closure in fluid simulations of collisionless reconnection}\citep{L256}

\textsf{本轮状态：}已有全文记录复用，不增加计数。

\textsf{核读范围：}期刊PDF1-16全部\par\textsf{图式：}图1-10及表1；图1、3-6、8-10与表1视觉复核

\href{run:../../papers/L256.pdf}{本地原文PDF}\quad\href{https://doi.org/10.1017/s002237781800048x}{DOI/版本入口}

\paragraph{实际模型与问题}全Maxwell十矩、对压力偏离各向同性的部分作Laplacian闭合，不是直接扩散完整温度；以Vlasov热流高阶矩作对照。

\paragraph{参数、初边值与实现}GEM质量比25、256×128基准及1024×512加密；另流体质量比400/1836。k0s分别为GEM3/ds、100×50di WHBG的1/(3ds)、岛合并1/(2ds)；不是一组参数全场景。岛lambda5/10/15/25di，两套网格512×256到2048×1024。

\paragraph{求解得到什么}原松弛闭合的热流散度形状和符号在某些区域错误；梯度闭合改善上游假热流与电子非对角压力。岛合并平均0-1.5tA率标度约lambda\textasciicircum{}(-.73)，出现反弹且无二次岛。表1lambda5至25平均率.180→.052；两套网格差小。GEM高分辨梯度流体后期斜率近Vlasov，但启动仍迟；高分辨局地松弛反而更慢。

\paragraph{成立条件与剩余缺口}k0按与动理学一致调节，尺寸标度也依赖k0；不能宣称无参预测。实际质量比只展示流体结构，不是同质量比完整动理学验证。显式Laplacian子循环中不交换块边界、不更新速度密度，作者承认块边界误差。GEM温度条件缺常规1/2，岛平衡坐标与边界文字不一致，复现需输入核对。

\paragraph{对Hall比较的意义}重要改进是直接检查热流矩再验证非线性反弹；其成功同时有问题相关的校准与数值限制。

\section{Ng2020：2020年}\label{paper:Ng2020}

\textbf{An improved ten-moment closure for reconnection and instabilities}\citep{Ng2020}

\textsf{本轮状态：}已有全文记录复用，不增加计数。

\textsf{核读范围：}1-8（作者稿全部，含附录A）\par\textsf{图式：}图1-7、式4-6、附录A；2026-11-14定向视觉复核全部结果图及附录页

\href{run:../../papers/Ng2020.pdf}{本地原文PDF}\quad\href{https://doi.org/10.1063/5.0012067}{DOI/版本入口}

\paragraph{实际模型与问题}十矩全Maxwell温度梯度热流，与局地各向同性化和非局地无磁化闭合比较；FV输运、局地隐式源、超时间步抛物子步。

\paragraph{参数、初边值与实现}Fadeev与Ng2015同系；lambda5di用640×320，同dx扫描系统尺寸，mi/me25, k0s=1/ds；导引场0-.4-1-2等。LHDI另例mi/me36, Ti/Te10, L=rhoi, vt, e=.06c, 盒6.4L×12.8L, 模数8。

\paragraph{求解得到什么}平均率随岛尺寸下降可接近PIC约lambda\textasciicircum{}-0.8；岛合并最大率随Bg升到1先升后降，不能套用Harris片单调下降。离子宽层及非旋转对称度改善。相同k0=1/ds可同时激发重联与LHDI。

\paragraph{成立条件与剩余缺口}LHDI增长率默认1.33Ωci、调k为ky后.86，动力学1.1，局地闭合.34，非局地.84；仍非完全吻合。电子各向异性空间结构仍不准、跨场热流过大；未研究LHDI对重联的非线性作用。附录指出零热流、无碰撞的对称稳态停滞点会驱动不物理Pxx趋零。 图6的流体电场图为显示除以2，不能仅由色彩图形近似声称幅度吻合。

\paragraph{对Hall比较的意义}梯度闭合并非更高阶就自动正确；改善来自k\textsuperscript{2}/|k0|的尺度选择，仍有自由参数和磁化缺失。

\section{TenBarge2019MMS：2019年}\label{paper:TenBarge2019MMS}

\textbf{An Extended MHD Study of the 16 October 2015 MMS Diffusion Region Crossing}\citep{TenBarge2019MMS}

\textsf{本轮状态：}已有全文记录复用，不增加计数。

\textsf{核读范围：}1-17（全部）\par\textsf{图式：}PDF5、6、9、10、11页直接视觉核验式12、图1/6-10；17页全文及13图图注读完。正式HTML新增第14图未视觉核验。

\href{run:../../papers/TenBarge2019MMS.pdf}{本地原文PDF}\quad\href{https://doi.org/10.1029/2019JA026731}{DOI/版本入口}

\paragraph{实际模型与问题}两物种全Maxwell五矩/十矩Gkeyll，局地Hammett–Perkins型压力松弛；直接有限矩流体证据。

\paragraph{参数、初边值与实现}mi/me100, c/vAs500；2D40.96×20.48di, 2048×1024；3D40.96×20.48×10.24di, 1024×512×256；周期双片；Bg/Bsx.099, Bmx/Bsx-1.7, Tem/Tes12.3, Tim/Tis5.63；k0s ds=.1，种子.1，磁噪声.0005。

\paragraph{求解得到什么}图1直接给离子/电子热能和流能，十矩2D/3D全域预算相近、离子热增量较大；电子各向异性峰约2.5近MMS幅度但位置在磁鞘侧而观测在磁层侧；局地闭合抑制LHDI，剪切/kink仍出现。

\paragraph{成立条件与剩余缺口}热能增量不是物种熵产Qi/Qe；五矩过快LHDI并明显总能损失，不可信作分配预测；未给网格/质量比独立收敛；nm/ns=16.7与低密度磁层文字冲突，l0=2π/k0≈d与k0d=.1冲突；不静默修正输入。

\paragraph{对Hall比较的意义}模型身份、温度/能量/熵产区别及跨文献判断见独立加热专题；不把本篇单独作为普适分配定律。

\section{E552：2021年}\label{paper:E552}

\textbf{Fluid simulations of three-dimensional reconnection that capture the lower-hybrid drift instability}\citep{E552}

\textsf{本轮状态：}已有全文记录复用，不增加计数。

\textsf{核读范围：}arXiv2008.06440v3 1-20全部，对应2021年期刊\par\textsf{图式：}图1-9；闭合式、图4、6-9视觉复核

\href{run:../../papers/E552.pdf}{本地原文PDF}\quad\href{https://doi.org/10.1017/s0022377820001683}{DOI/版本入口}

\paragraph{实际模型与问题}全Maxwell五矩与十矩比较；压力偏离各向同性的梯度闭合chi3, k0s1/ds；Vlasov用于二维电流片不稳定性对照。

\paragraph{参数、初边值与实现}CWENO+RK3，Maxwell FDTD，Laplacian二阶中心，每子循环交换压力边界，三维用8次；压力对角下限。三维100×25×6di、2400×600×150、mi/me100, c/VA30，Harris a.5, nb.2, Ti/Te5；磁通种子.1、电子密度噪1e-8。独立静电LHDI nb.001、Ti/Te10、a约.67、质量比36/100/250至流体1836。

\paragraph{求解得到什么}静电LHDI质量比100时Vlasov/十矩梯度/五矩/十矩松弛增长率2.2/2.4/3.2/.67 Omega\_i；最快模随质量比变化，固定单一模比较会误导。电磁支Vlasov约.41、流体.21，均出现超过B0的磁扰动及强kink。三维较二维启动早、后期片变宽，出流湍动强而电子层相对适中；密度噪1e-10至1e-2使LHDI启动约30→4 Omega\_i\textasciicircum{}-1。

\paragraph{成立条件与剩余缺口}高质量比无同参数Vlasov验证；电磁支结构、启动与增长仍不同，异常高饱和原因未解释。CSSI只是形态相似的候选，未确诊。三维率是z平均电场代理；电子de约2.4格，未给全三维网格收敛。对角下限不等于完整压力正定保证；印刷Lap约+k0\textsuperscript{2}应理解为负号Fourier符号。

\paragraph{对Hall比较的意义}该工作同时检验重联与不稳定性并强调噪声/启动竞争，优于仅比较峰值；但不能把3D形态相近称为完整动理学等价。

\section{Closure2026：2026年}\label{paper:Closure2026}

\textbf{Capturing Secondary Kinetic Instabilities in Three‐Dimensional Dayside Reconnection Using an Improved Gradient‐Based Closure}\citep{Closure2026}

\textsf{本轮状态：}已有全文记录复用，不增加计数。

\textsf{核读范围：}1-23（arXiv2602.16960v1全部，含附录A）\par\textsf{图式：}图1-11、附录A1及式A1-A12；2026-11-14定向视觉复核全部结果图及附录页

\href{run:../../papers/Closure2026.pdf}{本地原文PDF}\quad\href{https://doi.org/10.1029/2026ja035250}{DOI/版本入口}

\paragraph{实际模型与问题}Gkeyll三维非对称十矩；温度梯度闭合与局地闭合对照；顶点对称热扩散和alpha=.75限制器。

\paragraph{参数、初边值与实现}盒40.96×20.48×10.24di，全周期双电流片，1152×576×288，mi/me100, c/vA200, Bg/Bsx=.099；nm/ns=.062, Bxm/Bxs=-1.696, Tem/Tes12.3, Tim/Tis5.63；k0αdα约1，磁噪声.0005，psi种子.1。

\paragraph{求解得到什么}梯度闭合出现约5Ωci\textasciicircum{}-1内LHDI、约20后kink与通量绳/次级岛，局地同参不出现；三维率先峰后降，压力张量支撑X点非理想场。梯度闭合成本约3倍。 图6右（PDF11页）中，梯度闭合相对局地闭合的离子压力能及流能变化更小，电子压力能增长较快；这不等于已测得更大不可逆电子熵产。

\paragraph{成立条件与剩余缺口}温度各向异性峰值侧别与MMS/PIC不符，非旋转对称度与各向异性区域过宽，作者认为背景跨场扩散产生不物理各向异性。独立2D Vlasov对照kink仅数个离子周期出现，流体需数十周期，且遗漏小尺度模。正文未给质量比/网格独立收敛序列。

\paragraph{对Hall比较的意义}证明某些先前受抑制的横向不稳定性可在十矩中出现，不是准确恢复全部三维动理学。建议的磁场方向热流及弱各向同性化仍为未来改进。

\section{L244：2018年}\label{paper:L244}

\textbf{A reduced Landau-gyrofluid model for magnetic reconnection driven by electron inertia}\citep{L244}

\textsf{本轮状态：}已有全文记录复用，不增加计数。

\textsf{核读范围：}HAL02049658v1封面及正文1–26全读；物理页1–27\par\textsf{图式：}式2.1–4.21、Hamiltonian/Casimir推导及图1–7；阻尼算子、图3–6视觉复核

\href{run:../../papers/L244.pdf}{本地原文PDF}\quad\href{https://doi.org/10.1017/S002237781800051X}{DOI/版本入口}

\paragraph{实际模型与问题}从Scott型gyrofluid按epsilon低频、delta\textsuperscript{2}=me/mi、beta\_e\textasciitilde{}delta\textsuperscript{2}约化为六场；含离子FLR、电子双温、电子惯性和非局地Landau闭合。实际重联只算冷离子、Ni=Tperpi=0、Tperpe恒零的子模型。

\paragraph{参数、初边值与实现}delta=.1, beta\_e=.05, lambda3；线性半盒Lx46, Ly24pi, Lz96pi；非线性ky1/12, kz1/48，Ly12pi, Lz48pi。显式AB3、y/z并行；非线性热流刚性使步长至多缩小100倍，改用单螺度坐标y*=y+(kz/ky)z的等价二维推进。全文未列具体格点数及绝对步长。

\paragraph{求解得到什么}Landau算子用局地平行梯度和磁场方向二阶空间平均的Fourier分母，非简单固定k0扩散。阻尼降低平行温度扰动、线性增长率，高m斜模可稳定；同一gamma*t上岛面积曲线相似，意味着真实时间中发展变慢，而非未受影响。饱和岛大小近似不变，细密度/电流/温度结构变平滑。闭合去除自由能，未出现动理学参考中的重联—加热延迟。

\paragraph{成立条件与剩余缺口}斜单螺度不是多螺度三维湍流；没有热离子完整六场非线性验证。无Landau时数值能损约总能量.2\%，却约为同期磁能释放20\%，不可忽略。未出现延迟可能是闭合过度耗散，作者提出高阶闭合问题。自由能损失解释为未解析相混合/电子加热，非直接分布函数测量。

\paragraph{对Hall比较的意义}给出可推导、可检验的闭合成功与失败并存例子；必须同时检验时序、饱和、能量预算，不能以一项吻合宣称动理学等价。

\section{F635：2026年}\label{paper:F635}

\textbf{Counter-streaming heat-flux closure for electron-only collisionless magnetic reconnection}\citep{F635}

\textsf{本轮状态：}已有全文记录复用，不增加计数。

\textsf{核读范围：}arXiv2607.07106v1全部10页；公开三个输入文件已下载核对关键设置\par\textsf{图式：}图1-6全部视觉核查；特别核对图4最佳时间指标、图5独立色标

\href{run:../../papers/F635.pdf}{本地原文PDF}\quad\href{https://doi.org/10.48550/arxiv.2607.07106}{DOI/版本入口}

\paragraph{实际模型与问题}2026预印本：固定离子背景，两个电子十矩流体；对照单电子十矩及OSIRIS PIC。不是通常电子—离子双流体。

\paragraph{参数、初边值与实现}48de×64de双周期，480×640；力自由双层半厚de，Brec=.5、Bg/Brec=.1、beta\_rec=.1。PIC每格625粒子、二阶形函数、dt=.066667/omega\_pe；流体Gkeyll分维有限体积CFL2/3、k=.1/de各向同性松弛。两电子群按片上下来源标记，tanh分配宽de、密度下限.03n0；初始速度相同。Zenodo输入中plasmaBetae0=.3等旧声明不参与实际温度初始化；实际vth=.111803与Brec=.5给beta=.1。

\paragraph{求解得到什么}两群相对漂移与不同压力可重建总三阶中心矩的大部分结构，改善单流体的密度空洞与局部压力堆积；热流符号/形状更接近PIC。在示例中单群内部热流对局部Qxxx贡献约4\%；但双流体重建热流幅度仍偏高。

\paragraph{成立条件与剩余缺口}图4是每个PIC参数算例在时间内取最大相似度，非双流体全参数性能或时间平均；图5流体色标不同，不能直接认定幅度一致。高beta时人为分群梯度引发自由膨胀，晚期同群再度多束化；三维、动态分裂和离子耦合均未验证。156倍是按粒子相空间变量数与20矩的表示压缩估算，不是实测运行加速。仍有弱松弛而非严格Q=0。

\paragraph{对Hall比较的意义}有价值的新方向是利用来源分群减少每群的闭合负担；不能推成单电子流体在数学上不可能表达这些矩，或热流闭合已解决。

\section{Wetherton2019：2019年}\label{paper:Wetherton2019}

\textbf{Validation of Anisotropic Electron Fluid Closure Through In Situ Spacecraft Observations of Magnetic Reconnection}\citep{Wetherton2019}

\textsf{本轮状态：}新增全文核读。

\textsf{核读范围：}OSTI接收稿PDF1–20全部；封面、正文、方法、讨论、参考文献、后部重复图页均读完，无文内附录。\par\textsf{图式：}图1–3、式1–6；直接查看PDF4、6、8、13、15、17–20页，图2密度转折与图3截面特别核对。

\href{run:../../papers/Wetherton2019.pdf}{本地原文PDF}\quad\href{https://doi.org/10.1029/2019GL083119}{DOI/版本入口}

\paragraph{实际模型与问题}各向异性电子EoS的MMS观测检验；H3D数值部分为动理学离子加流体电子，非完整两流体。外部来源、俘获/通过粒子决定p\_parallel(n, B)与p\_perp(n, B)，未演化全非旋转对称压力。

\paragraph{参数、初边值与实现}2015-10-31强导场近对称MMS3事件，轨迹距X点约7–10di。来源n=7.1cm\textasciicircum{}-3、B=55.5nT、Tperp=32.4eV、Tparallel/Tperp=1.22。H3D力自由初态，Bg/Br=1.75、总场beta\_e=.02、Ti/Te=9.1；80×40di、2048×1024网格、400离子粒子/格；显式电阻破冻结。t=80/Omega\_ci选宽约3.5di截面对照；本文未列完整边界、时间步、电阻值或收敛序列，需前序方法与输入。

\paragraph{求解得到什么}观测平行温度在低密度近Boltzmann、高密度近CGL，EoS跨越转折；混合模拟避免低质量比100 PIC在密度空洞处的过强平行降温。强导场力平衡加绝热离子给归一化电子温升弱beta依赖，典型条件约.02；这是温升而非两物种不可逆热量比。

\paragraph{成立条件与剩余缺口}数值验证属混合，不能计为完整双流体。来源与截面匹配使其不等同无参数预测；未给统一残差/仪器协方差。峰值B≈Bg、Ti\ensuremath{\gg}Te及绝热离子限制标度；图2c横轴电子重联场beta，式4–6中beta\_r为离子beta，不能混用。式1–2周围平行热流符号与三阶矩定义含二分之一规范疑点，报告采用q\_a=Q\_abb/2。前序EoS推导、原始MMS数据和方法论文未完整重读。

\paragraph{对Hall比较的意义}强于单峰温度符合的证据是跨密度区间的各向异性转折及分布边界。支持电子闭合，不能验证离子热化或通用分配。相对标量Hall新增的是压力自由度与受俘获限制的热流响应；可以嵌入各向异性Hall，非对所有Hall的否定。

\section{Wetherton2021：2021年}\label{paper:Wetherton2021}

\textbf{Anisotropic Electron Fluid Closure Validated by in Situ Spacecraft Observations in the far Exhaust of Guide-field Reconnection}\citep{Wetherton2021}

\textsf{本轮状态：}新增全文核读。

\textsf{核读范围：}OSTI接收稿PDF1–23全部；正文、模型、全部图表、讨论、参考文献读完，无文内附录。\par\textsf{图式：}图1–10与表1、式1–4；直接看PDF6–7、9、11–12、14–19页，重点图6–7的偏差区和图9–10的反演增能。

\href{run:../../papers/Wetherton2021.pdf}{本地原文PDF}\quad\href{https://doi.org/10.1029/2020JA028604}{DOI/版本入口}

\paragraph{实际模型与问题}MMS远下游观测与电子各向异性EoS比较；数值取矩求压力而非完整双流体重联计算。论文第2节复用PIC解释密度与电势，不能当作新增纯流体模拟。

\paragraph{参数、初边值与实现}主事件2016-01-21约01:06UT，MMS2，距X线约100di、Bg/Br约.7、略非对称；以2015-10-31近区事件为对照。分别选择入流1/出流/入流2来源n=47/41/40cm\textasciicircum{}-3，B=30/29/41nT，Tperp=28/30/30eV，Tparallel/Tperp=.9/1/.9以匹配温度。各区域数百采样，但四个来源参数不是独立直接远端测量。复用PIC图1质量比1836、20×20di；本文不提供新流体网格收敛。

\paragraph{求解得到什么}三来源EoS大体解释平行与垂直温度随局地n、B的变化，给出三条不同区域特征曲线。部分垂直温度明显超过预测，作者联系到区域磁压缩、局地俘获和有限通过时间，以DeltaE=mu积分partial\_t B dt反演额外增能；深俘获例1.35与B/Binf≈1.34相符，出流额外项通常小于.1。

\paragraph{成立条件与剩余缺口}来源参数按观测选取；反演DeltaE亦使用观测，不能视作独立预测。原文明确DeltaE没有对n、B的形式函数，未形成通用新EoS。压缩起源、是否延伸至X线不明；局地磁场极小也不必总有大增能。未给通用拟合残差、置信区间或完整诊断响应。图7b曲线有0/10/20eV人为偏移，不可读为绝对温差。前序时序定位和仪器论文、原始MMS数据未复算。

\paragraph{对Hall比较的意义}把近X线闭合支持扩展到远区同时给出失败边界。不能把经三来源拟合的吻合当成任意条件下预言；额外垂直增温需要历史或新增状态信息。对标量Hall是各向异性约束，对既有EoS-Hall则是局地闭合与来源假设的限制；未解决离子电子不可逆热化比。

\section{Tanabe2015MAST：2015年}\label{paper:Tanabe2015MAST}

\textbf{Electron and Ion Heating Characteristics during Magnetic Reconnection in the MAST Spherical Tokamak}\citep{Tanabe2015MAST}

\textsf{本轮状态：}已有全文记录复用，不增加计数。

\textsf{核读范围：}所记录版本全部5个PDF物理页；正文、方法、结果、讨论、图表、参考文献及适用文内附录均核读。\par\textsf{图式：}直接查看PDF物理页2, 3, 4；All Figs1–5;  experimentgeometry, Te/ne/Titomography, timeevolution, guidefieldscan

\href{run:../../papers/Tanabe2015MAST.pdf}{本地原文PDF}\quad\href{https://doi.org/10.1103/PhysRevLett.115.215004}{DOI/版本入口}

\paragraph{实际模型与问题}本文是球形托卡马克的强导场双环合并实验，以汤姆孙散射测量电子温度和密度，以离子多普勒层析重建温度空间分布，并与已有双流体模拟比较。实验同时讨论电子与离子的碰撞能量交换、各向异性热输运及封闭磁面约束。论文自身不进行独立霍尔模拟，温度结构的符合是对具体模型解释的实验约束。

\paragraph{参数、初边值与实现}环向导场约为 0.3 至 0.8 特斯拉，重联场约为 0.07 至 0.15 特斯拉；离子皮肤深度约 0.1 米，离子回旋半径小于 0.01 米。电子诊断分别为 300 和 130 通道，空间分辨率约 15 与 10 毫米；离子诊断为 32 弦碳离子光谱层析。二维剖面组合多个放电脉冲和测量位置，重联场依据合并后磁平衡重建估计。

\paragraph{求解得到什么}电子在重联中心附近形成宽约 2 至 5 厘米的热点，温度在小于 1 毫秒内由约 10 升到 200 电子伏；离子温升分布在出流及下游停滞区。经过约 4 至 10 毫秒的两物种能量交换时间，两个温度剖面均形成三峰。更高导场使中心电子热点更尖，但所测下游离子温度基本不变，与较弱导场机制的趋势不同。

\paragraph{成立条件与剩余缺口}原文明确指出缺乏直接磁探针测量限制了热点成因的判别。黏性、压缩阻尼与输运解释依赖模型，温度图并不是完整局地功率或熵产预算；杂质离子温度与主离子温度的对应也依赖诊断和碰撞条件。全部二维点并非单次放电同时测得；独立诊断论文、学位论文和原始脉冲数据未在本次重新分析。

\paragraph{对Hall比较的意义}这是双温流体建模线的重要实验支撑：原文直接引用 L223，并将下游温升与含碰撞黏性耗散的双流体结果联系。但热点位置相似不足以唯一验证闭合，强导场下约束增强和驻留时间延长也会改变最终温度。该文支持具体装置中的局部加热研究，没有给出普适物种分配比；正式版 5 页及全部五幅图已核读。

\section{F674：2015年}\label{paper:F674}

\textbf{Study of energy conversion and partitioning in the magnetic reconnection layer of a laboratory plasma}\citep{F674}

\textsf{本轮状态：}已有全文记录复用，不增加计数。

\textsf{核读范围：}所记录版本全部11个PDF物理页；正文、方法、结果、讨论、图表、参考文献及适用文内附录均核读。\par\textsf{图式：}直接查看PDF物理页6, 7, 8, 9；Figs6–8,  TableI,  Eqs(2)–(6)

\href{run:../../papers/F674_published.pdf}{本地原文PDF}\quad\href{https://doi.org/10.1063/1.4920960}{DOI/版本入口}

\paragraph{实际模型与问题}本文通过磁重联实验装置的局部测量建立有限控制体能量收支，并与装置几何下及开放边界下的全动理学粒子模拟比较。分别积分电子和离子的电场功，再区分有序流能、热储能与边界输运。它是实验能量预算和动理学对照研究，不能列为霍尔流体独立预测温度比例的数值证据。

\paragraph{参数、初边值与实现}本次读取 2015 年正式版 11 页，保留的 2014 年预稿按同一论文处理。开放粒子模拟采用真实质子质量比 1836，初始离子与电子温度比为 5，背景密度比为 0.23，背景电子温度比为 0.76。模拟盒为 20 乘 20 个参考离子皮肤深度，网格为 5120 乘 5120，每种粒子约一百亿个；电流片半宽为 0.5 个参考离子皮肤深度。

\paragraph{求解得到什么}图 7 给出离子和电子电场沉积率分别为 0.71±0.15 与 0.35±0.06 兆瓦，约为二比一；两物种的热能变化项却分别约为 0.27±0.10 与 0.26±0.10 兆瓦。后者还包含热储能变化和边界焓流差。表一改变控制体后，离子和电子占输入磁能的份额也发生变化，说明该比例与统计区域有关。

\paragraph{成立条件与剩余缺口}预算中的热损失具有较大误差，部分离子损失只给出上限；输出能流、边界条件和采样尺度都会影响份额。粒子模拟虽然包含真实质量比，也未覆盖任意导场、等离子体压强、三维湍流及宏观装置。本文承接 2014 年实验成果，不能把两篇文章当成相互独立的复现实验；原始数据和模拟代码未在本次复算。

\paragraph{对Hall比较的意义}常被引用的二比一是物种电场能量沉积率比，既不是离子与电子温度比，也不是两物种不可逆熵产生热比。本文证明在指定实验及控制体内可以建立有用的分物种能量预算，但没有提供普适加热闭合。应同时报告有序流动、热储能、焓流和热损失，不能只摘取一个比值作为霍尔模型的普遍预测。

\section{F679：2014年}\label{paper:F679}

\textbf{Bulk ion acceleration and particle heating during magnetic reconnection in a laboratory plasma}\citep{F679}

\textsf{本轮状态：}已有全文记录复用，不增加计数。

\textsf{核读范围：}所记录版本全部13个PDF物理页；正文、方法、结果、讨论、图表、参考文献及适用文内附录均核读。\par\textsf{图式：}直接查看PDF物理页5, 8, 9, 10, 11；Figs 5,  9–13;  Eqs (1)–(19)

\href{run:../../papers/F679.pdf}{本地原文PDF}\quad\href{https://doi.org/10.1063/1.4874331}{DOI/版本入口}

\paragraph{实际模型与问题}本文是磁重联实验装置的氦等离子体实验，并以二维全动理学粒子模拟进行对照。电子动量平衡用于解释面内电势，光谱与探针分别约束离子温度和流速；模拟比较有、无离子碰撞的速度分布。文中的“双流体区域”指电子与离子解耦的实验状态，并不意味着所用数值模型是有限矩双流体模型。

\paragraph{参数、初边值与实现}实验气压为 4.5 毫托，磁通芯间距 42 厘米，离子皮肤深度约 8 厘米，无外加导场。4200 次放电约三成通过筛选，每点平均 7 至 15 次；离子光谱的空间、时间分辨率约为 3 至 4 厘米和 5.6 微秒。粒子模拟质量比为 400，初始两物种等温，电子热速为光速的 0.183 倍；离子碰撞频率及自由程相对特征尺度与实验匹配。

\paragraph{求解得到什么}面内电场约 700 伏每米，高于约 200 伏每米的重联电场，主要离子增能由面内电势驱动。有序流能小于电势能下降量，下游出现更高离子温度。无碰撞模拟的较大温度包含多束混合；加入实验量级碰撞后，分布更接近麦克斯韦形态。电子能量沉积约 340 千瓦，经典电阻贡献估计仅约 59 千瓦。

\paragraph{成立条件与剩余缺口}实验没有完整测得电子压力张量及热通量，约 300 千瓦的热损失估计具有较大误差，因此异常电子耗散是条件性推断，不能唯一归因于所见波动。二维模拟的部分反向离子束也未被实验空间分辨率辨认。完整网格、时间步和边界细节部分依赖前序方法文献，本次未复核全部前序实现或复算原始数据。

\paragraph{对Hall比较的意义}本文把电势加速与下游分布热化联系起来，支持轨道运动、再磁化和碰撞共同作用的解释，为流体闭合提供实验约束。它没有证明只加入霍尔项就能得到正确熵产，也没有给出普适离子与电子不可逆加热比。应把物种电场功、流动能、压力功和热输运分开核算；本次依据正式版全部 13 页及关键图式作出判断。

\section{Yoo2017Heating：2017年}\label{paper:Yoo2017Heating}

\textbf{Electron heating and energy inventory during asymmetric reconnection in a laboratory plasma}\citep{Yoo2017Heating}

\textsf{本轮状态：}已有全文记录复用，不增加计数。

\textsf{核读范围：}所记录版本全部18个PDF物理页；正文、方法、结果、讨论、图表、参考文献及适用文内附录均核读。\par\textsf{图式：}直接查看PDF物理页7, 8, 10, 12, 14, 15；Figs4, 5, 8, 10, 11;  Table1;  Eqs5–10；2026-11-13纠正定位：本地PDF图10、11在第15页，式(10)在第14页，图8在第12页。

\href{run:../../papers/Yoo2017Heating.pdf}{本地原文PDF}\quad\href{https://doi.org/10.1002/2017JA024152}{DOI/版本入口}

\paragraph{实际模型与问题}本文研究两侧密度不同的磁重联实验，并使用二维和三维全动理学粒子模拟解释电子加热及物种能量收支。实验重点测量低密度侧分离面附近的电场、温度和波动；模拟比较温度矩与速度分布形状的变化。因此它既约束空间非对称条件下的能量分配，也检验分布趋于麦克斯韦形态是否必然伴随额外净增能。

\paragraph{参数、初边值与实现}实验两侧密度比约为 8，重联磁场分别约为 145 与 168 高斯，残余导场小于重联场的一成，离子皮肤深度约 7.5 厘米。粒子模拟质量比为 900、密度比为 10，初始两物种等温，高密度侧热压与磁压比为 0.5。二维盒长为 10 乘 10 个离子皮肤深度、网格 1024 平方；三维第三方向长 3 个离子皮肤深度，增加 320 层网格。

\paragraph{求解得到什么}低密度侧分离面具有增强的电子加热。图 10 中电子和离子沉积率分别为 0.35±0.07 与 0.44±0.09 兆瓦，占输入磁能约百分之二十五和三十一；两物种热能变化项均约 0.30 兆瓦。二维与三维模拟温度剖面相近，三维电子分布却更接近麦克斯韦形态，说明分布热化与净能量增长可以表现不同。

\paragraph{成立条件与剩余缺口}实验波动与温升的相关系数约为 0.29，显著性水平约为 0.086，不能确立低混杂漂移波的强因果作用。模拟的电子温度、频率比和有限盒尺度与实验仍有差异。正文温度矩定义及表格密度单位存在排印疑点，应按标准定义和正文量级核对，不能据此反推代码错误。原始数据入口已记录，但本次未重新分析。

\paragraph{对Hall比较的意义}非对称实验下两物种沉积率相近，直接限制了对称实验约二比一结果的外推。三维波动可能促进分布散射而不显著增加温度，这提示流体闭合应分别接受能量预算和分布诊断检验。本文支持针对具体条件研究物种分配，但未建立跨装置、跨采样位置的普适比例；正式版 18 页及能量收支图已完整核读。

\section{Che2014Buneman：2014年}\label{paper:Che2014Buneman}

\textbf{Two-Fluid Description of Wave-Particle Interactions in Strong Buneman Turbulence}\citep{Che2014Buneman}

\textsf{本轮状态：}已有全文记录复用，不增加计数。

\textsf{核读范围：}1-26（全部）\par\textsf{图式：}PDF7、16-19、21页视觉核验图1、6-9和完整流体方程；所有26页含图注/参考文献读完，无附录。

\href{run:../../papers/Che2014Buneman.pdf}{本地原文PDF}\quad\href{https://doi.org/10.1063/1.4882677}{DOI/版本入口}

\paragraph{实际模型与问题}3D PIC诊断与近似双流体闭合分析；未独立演化验证纯双流体重联，模拟未发展重联。

\paragraph{参数、初边值与实现}mi/me100, w0=.1di=de, Te=Ti=.04miVA\textsuperscript{2}；盒1×1×2di，512×512×1024；x/z周期、y导电；电子漂移9VA≈3vte、离子.9VA；|B|≈sqrt26 B0。

\paragraph{求解得到什么}Buneman拖曳和相混把电子漂移能转为平行热，磁能下降伴电子垂直热增长；局地流体矩预算近似刻画PIC；离子主要协助波交换，弱热化。

\paragraph{成立条件与剩余缺口}Dez(t)仍需动力学理论/PIC拟合；同前序Che2013运行，非独立验证；无实际重联、无实验比；式19为电子平行/垂直比。图7约化PIC熵增不等于精细Vlasov熵增，未熵粗粒化收敛验证。

\paragraph{对Hall比较的意义}模型身份、温度/能量/熵产区别及跨文献判断见独立加热专题；不把本篇单独作为普适分配定律。

\chapter{访问、版本及未核读线索}

Ohia2012的MIT公开记录已取得，PDF入口返回HTTP 429；OSTI旧入口未返回匹配PDF。Ohia2015出版PDF入口返回403，网页只用于检索和追引，尚未完成原文图式核验；二者不计新全文，不作为本专题科学结论的直接证据。其模型能力仍不能仅由后继论文转述代替原文核读。

Kaptanoglu2020勘误与TenBarge2019MMS正式PDF本轮访问仍未成功，沿用已读接收稿或arXiv版及已有正式网页差异记录。十三/二十矩相关会议摘要、Ohia学位论文和前序EoS原始推导保留为追引；不声称已读、已验证或已复算。

Wetherton2019使用20页OSTI接收稿，后部含重复/放大图页；Wetherton2021使用23页OSTI接收稿。两篇均无文内附录。独立仪器论文、原始MMS数据及前序方法代码未在本轮完整重新分析。
\backmatter
\bibliographystyle{unsrtnat}
\bibliography{references}
\end{document}
